\documentclass[11pt,onecolumn]{IEEEtran}
\usepackage{amsmath,amssymb,amsthm}
\usepackage[hidelinks]{hyperref}

\newtheorem{theorem}{Theorem}
\newtheorem{proposition}[theorem]{Proposition}
\newtheorem{corollary}[theorem]{Corollary}
\newtheorem{lemma}[theorem]{Lemma}
\newtheorem{remark}[theorem]{Remark}

\newcommand{\E}{\mathbb{E}}
\newcommand{\R}{\mathbb{R}}
\newcommand{\C}{\mathbb{C}}
\newcommand{\Z}{\mathbb{Z}}
\newcommand{\tr}{\operatorname{tr}}
\newcommand{\Cov}{\operatorname{Cov}}
\newcommand{\hY}{\hat{\mathbf Y}}
\newcommand{\dw}{\tfrac{d\omega}{2\pi}}
\newcommand{\VG}{V_{\mathrm G}}
\newcommand{\bA}{\mathbf A}
\newcommand{\bC}{\mathbf C}
\newcommand{\bE}{\mathbf E}
\newcommand{\rhid}{\rho_{\mathrm{hid}}}

\emergencystretch=1.5em
\begin{document}

\title{Rate, Leakage, and Distortion for Unstable and Rapidly Varying Gaussian Sources}

\author{Andrew~H.~Bond,~\IEEEmembership{Senior Member,~IEEE}%
\thanks{A. H. Bond is with the Department of Computer Engineering, San Jos\'e
State University, San Jos\'e, CA 95192 USA (e-mail: andrew.bond@sjsu.edu).
This is a working draft; the title is provisional.}}

\markboth{}{Bond: Rate, Leakage, and Distortion for Unstable and Rapidly Varying Gaussian Sources}

\maketitle

\begin{abstract}
An encoder describes a Gaussian source to a decoder under a weighted
quadratic distortion, while a second party holds a noisy observation of
the source. The description has two costs, its rate
and its conditional leakage to the second party. The region of
achievable pairs is known for stationary sources, for sources whose
spectrum varies slowly or periodically, and for scalar unstable
autoregressive sources. We treat two open cases. For a vector
autoregressive source with unstable poles, observed by the second party
through a linear filter of finite memory in white noise, the region is that of a
stationary source moved by two amounts. In rate it moves by the sum of
the logarithms of the moduli of all unstable poles, and in leakage by
the same sum over only the poles that the second party cannot observe.
Instability is therefore free in leakage exactly when every unstable
pole is observable. The proof writes a window of the source as a
stationary bridge together with its final state. For a scalar source
whose window covariances have a limiting eigenvalue distribution,
observed through a scalar gain in white noise, the region is a single
quadrant given by reverse water-filling on that distribution. This
covers statistics that vary from symbol to symbol without a period.
For a cyclostationary process sampled at a rate incommensurate
with its period, the limit exists and does not depend on the sampling
phase, which settles a published conjecture.
\end{abstract}

\begin{IEEEkeywords}
Rate--distortion, information leakage, side information, unstable
sources, Kalman filtering, almost periodic processes, cyclostationary
processes.
\end{IEEEkeywords}

\section{Introduction}
\label{sec:intro}

This paper continues \cite{part1,part2,part3}. An encoder sees a source
and sends a message to a decoder, which reproduces the source under a
weighted quadratic distortion. A third party, the holder, observes the
source through a linear map in Gaussian noise. It does not decode, but
the message tells it something, and the leakage is the information the
message gives about the source beyond what the holder already has. The
operational region of pairs of rate and leakage was found in
\cite{part1} for Gaussian vectors, in \cite{part2} for stationary
Gaussian processes, and in \cite{part3} for sources whose spectrum varies
slowly in time, for cyclostationary sources, and for scalar unstable
autoregressive sources. The last paper ended with a list of open
problems. This one solves two of them.

The first is the unstable vector source. A scalar autoregressive source
with poles outside the unit circle costs, in rate, the logarithms of the
moduli of those poles on top of the stationary cost, and in leakage it
costs the same amount only when the holder does not observe the source
at all \cite[Theorem~26]{part3}. That proof separated the problem over
the eigenvalues of the window covariance, which works only because a
scalar source, a scalar weight, and a scalar holder commute. For a
vector source they do not. We show that the region is still a translate
of a stationary region (Theorem~\ref{thm:vector}). The rate moves by
$\rho_\star$, the sum of $\log|\rho|$ over the unstable poles $\rho$,
and the leakage moves by $\rhid$, the same sum over the unstable poles
that are hidden from the holder in the sense of Kalman: their modes do
not reach its observation. The holder may observe the source through
any filter of finite memory, and the filter can hide a pole by
cancelling it; a holder that sees every innovation, in noise, learns
nothing about the growth. An unstable pole is paid for once, in rate,
when the holder can see it, and twice when it cannot. The amount $\rhid$
jumps when the holder's filter is perturbed so that a hidden mode
becomes visible, however weakly, because the visible part of a growing
mode eventually rises above any noise.

The proof rests on a structural fact that does not need the eigenvalues
to separate. The inverse covariance of a window of the source differs
from a block Toeplitz matrix only in its last few blocks
(Lemma~\ref{lem:bridge}). So the source in the window, given its final
state, has the law of a stationary bridge, and all of the growth is
carried by the final state. The cost of describing that state is its
growth in volume, $\rho_\star$ per symbol, and the information it leaks
is its growth in volume as seen by the holder, $\rhid$ per symbol
(Lemma~\ref{lem:endpoint}). The second number comes from an identity
for the innovations of the Kalman filter (Lemma~\ref{lem:innovations}).
The window cost itself is compared with that of a block circulant
problem, in which the source and the holder commute with a cyclic
shift and the problem splits over frequencies
(Lemmas~\ref{lem:perturb} and~\ref{lem:circulant}); the two differ by
the two growth rates times $N$, up to $o(N)$.

The second problem is fast variation without a period. Between a
spectrum that drifts slowly and one that repeats are sources whose
statistics change from one symbol to the next and never repeat. A
cyclostationary signal sampled at a rate incommensurate with its period
is the standard example \cite{abakasanga2020,tan2024}. For a scalar
source and a scalar holder in white noise we show that the region
depends only on the limiting distribution of the eigenvalues of the
window covariances, when it exists, and is a single quadrant
(Theorem~\ref{thm:law}). For sources driven by a uniquely ergodic map,
such as the rotation that incommensurate sampling produces, that limit
exists and does not depend on the starting point
(Proposition~\ref{prop:ergodic}). For the sampled cyclostationary
sources of \cite{tan2024}, this proves the conjecture stated there that
the rate--distortion function does not depend on the sampling phase,
and it gives the function in closed form (Corollary~\ref{cor:sampled}).
For a vector source and a matrix holder the cost no longer depends on
eigenvalues alone. There the region is the limit of the costs of
independent blocks whose phases follow the invariant measure, again the
same for every starting phase (Theorem~\ref{thm:ergodic}); the error of
blocks of length $n$ is $O(1/n)$.

\textbf{Related work.} The rate--distortion function of a scalar
Gaussian autoregressive source with unstable poles is due to Gray
\cite{gray1970} and to Hashimoto and Arimoto \cite{hashimoto1980}, and
Gray and Hashimoto showed that their two forms agree
\cite{grayhashimoto2008}. Sahai and Mitter gave an upper bound for
first-order sources and sketched the vector case
\cite{sahai2007}. Derpich, M\"uller, and \O stergaard extended the scalar
result to other driving noises through the entropy gain of linear
systems \cite{derpich2021}. Guti\'errez-Guti\'errez and coauthors found
the rate--distortion function of vector sources that are asymptotically
wide-sense stationary \cite{gutierrez2018}, a class that excludes
unstable poles, and computed it for nonstationary vector ARMA sources
whose companion matrix has spectral radius below one
\cite[Theorem~4]{gutierrez2024}, which excludes them as well. Toms and
Berger gave the rate--distortion function of linear, finite-memory
Gaussian dynamic systems and showed that at small distortions it equals
that of the driving noise \cite{toms1971}. We have read only the
abstract of that correspondence and do not know whether its systems may
be unstable; at small distortions the rate of Theorem~\ref{thm:vector}
agrees with that statement. We did not find the rate--distortion
function of an unstable vector autoregressive source at every
distortion stated elsewhere, and we know of no result on leakage to an
observer for any unstable source. The identity of
Lemma~\ref{lem:innovations} is a discrete-time relative of the limits
of Braslavsky, Seron, Mayne, and Kokotovi\'c for continuous-time
Kalman--Bucy filters \cite{braslavsky1999}. For sampled cyclostationary
sources, Abakasanga, Shlezinger, and Dabora found the rate--distortion
function of memoryless samples at small distortions as a limit
superior of synchronous approximations \cite{abakasanga2020}, and Tan,
Dabora, and Poor extended it to sources with memory as an average over
the sampling phase of limits superior of finite programs, conjecturing
that the phase does not matter \cite{tan2024}.

Section~\ref{sec:model} fixes the setting and lists the results used.
Section~\ref{sec:vector} treats unstable vector sources and
Section~\ref{sec:law} sources with a limiting eigenvalue distribution.
Section~\ref{sec:open} lists what remains open.

\section{Setting and Imported Results}
\label{sec:model}

Logarithms are base two unless written $\ln$, and entropies are in
bits. Norms of matrices are operator norms. Notation follows
\cite{part1,part2,part3}. The decoder reproduces the source itself with
a weight $Q\succ0$, and $q>0$ denotes its smallest eigenvalue.

\subsection{Gaussian vectors of several symbols}

Let $T=(T_1,\dots,T_n)\in\R^{pn}$ be a centered Gaussian vector with
covariance $\mathsf S\succ0$, and let the holder observe $S=\mathsf GT+U$,
where $\mathsf G$ is any real matrix and $U\sim\mathcal N(0,I\otimes\Sigma_U)$,
$\Sigma_U\succ0$, is independent of $T$. Its information matrix is
$\mathsf J=\mathsf G^{\mathsf T}(I\otimes\Sigma_U^{-1})\mathsf G$. A code for
the vector has an encoder that maps $T$, and private randomness
independent of $T$ and of the noise, to a message $M$ in a finite set,
and a decoder that maps $M$ to reproductions $\hat T_t$. Its rate,
distortion, and leakage per symbol are $\tfrac1n\log|\mathcal M|$,
$\tfrac1n\sum_t\E\|T_t-\hat T_t\|_Q^2$, and $\tfrac1nH(M\mid S)$. For a
symmetric $\Delta$ with $0\prec\Delta\preceq\mathsf S$ the two costs are
\begin{equation}
r(\Delta)=\tfrac12\log\frac{\det\mathsf S}{\det\Delta},\qquad
\ell(\Delta)=\tfrac12\log\frac{\det(\Delta^{-1}+\mathsf J)}{\det(\mathsf S^{-1}+\mathsf J)},
\label{eq:costs}
\end{equation}
and for $D>0$ the frontier cost is
\[
\VG(\alpha,nD;\mathsf S,\mathsf J)=\min\bigl\{\alpha r(\Delta)+(1-\alpha)\ell(\Delta):\ 0\prec\Delta\preceq\mathsf S,\ \tr\bigl((I_n\otimes Q)\Delta\bigr)\le nD\bigr\}.
\]
This is the Gaussian vector problem of \cite{part1} with the holder
matrix $\mathsf G$ and the noise covariance $I\otimes\Sigma_U$. For a code,
$M-T-S$ is a Markov chain because the private randomness is independent
of the noise, and $n R\ge H(M)\ge I(T;M)$ and $nL=H(M\mid S)\ge I(T;M\mid S)$.
By \cite[Proposition~2]{part1}, applied to the variable $M$ from which
the reproductions are computed, every code for the vector, with or
without private randomness, therefore has
$\alpha R+(1-\alpha)L\ge\tfrac1n\VG(\alpha,nD;\mathsf S,\mathsf J)$ at its
distortion $D$, and by \cite[Theorem~5]{part1} both costs are convex
in $\Delta$ for any $\mathsf J\succeq0$.

\subsection{Stationary sources and a holder of finite memory}

The holder observes
\[
S_t=\sum_{j=0}^mH_jT_{t-j}+U_t,
\]
with $m\ge0$, real $r\times p$ matrices $H_j$, and $U_t$ independent
$\mathcal N(0,\Sigma_U)$; the case $m=0$ is a memoryless map. Write
$G(z)=\sum_jH_jz^j$ and $j(\omega)=G(e^{i\omega})^*\Sigma_U^{-1}G(e^{i\omega})$.
For a continuous spectral density $f$ on the circle with positive
definite $p\times p$ values, let $V_f(\alpha,D)$ and $\mathcal R_f(D)$ be
the spectral cost and spectral region of \cite{part2} for the
stationary Gaussian source with density $f$, the weight $Q$, and this
holder, with $D_{\max}^f=\int\tr(Qf)\,\dw$. The spectral cost is the
minimum of the integral over frequency of the costs \eqref{eq:costs} of
the Gaussian problem with covariance $f(\omega)$ and information
matrix $j(\omega)$, under a common distortion budget. It depends on
$(f,j)$ only through the joint distribution of their values, so the
orientation of the frequency axis does not matter. The function
$V_f(\alpha,\cdot)$ is convex, nonincreasing, continuous, and zero from
$D_{\max}^f$ on. We use the following results of \cite{part2}.
\begin{itemize}
\item[(i)] For $0<D<D_{\max}^f$ every point of $\mathcal R_f(D)$ is
approached by codes for the stationary source with deterministic
encoders and blocklengths tending to infinity, the holder observing
$S_t$ at the times of the window: \cite[Theorem~16]{part2} for a
memoryless map and \cite[Theorem~20]{part2} for a filter of finite
length. For $t\le m$ these observations involve symbols before the
window.
\item[(ii)] For probability measures $\mu_n\to\mu$ weakly on a compact
set carrying continuous data, the costs of the programs over $\mu_n$
converge to that over $\mu$ for $0<D<\int\tr(Qf)\,d\mu$
\cite[Lemma~2]{part2}. In particular, the program over the uniform
measure on the $N$ Fourier frequencies tends to $V_f$.
\item[(iii)] If $f'\preceq f\preceq(1+\kappa)f'$ at every frequency,
then
$V_{f'}(\alpha,\tfrac D{1+\kappa})-\tfrac p2\log(1+\kappa)\le V_f(\alpha,D)\le V_{f'}(\alpha,D)+\tfrac p2\log(1+\kappa)$.
This is the inequality of \cite[Lemma~11(c)]{part2}, and its proof is
two lines with the costs \eqref{eq:costs} at each frequency, for any
$j$. For the right inequality, the optimal field $\Delta$ for $f'$ is
feasible for $f$, and since $f^{-1}\preceq f'^{-1}\preceq(1+\kappa)f^{-1}$,
its rate and leakage for $f$ exceed those for $f'$ by at most
$\tfrac12\log\det(f'^{-1}f)\le\tfrac p2\log(1+\kappa)$ and
$\tfrac12\log\det\bigl((f'^{-1}+j)(f^{-1}+j)^{-1}\bigr)\le\tfrac p2\log(1+\kappa)$.
For the left one, the optimal field $\Delta$ for $f$ divided by
$1+\kappa$ is feasible for $f'$ at the distortion $D/(1+\kappa)$, and
its costs for $f'$ exceed those of $\Delta$ for $f$ by at most
$\tfrac p2\log(1+\kappa)$, because $f'\preceq f$ and
$(1+\kappa)\Delta^{-1}+j\preceq(1+\kappa)(\Delta^{-1}+j)$.
\end{itemize}
Achievable pairs and the closure $\mathcal A(D)$ of the achievable set
are defined as in \cite[Sec.~VI]{part3}, with codes of blocklengths
tending to infinity and private randomness allowed.

\section{Unstable Vector Autoregressive Sources}
\label{sec:vector}

\subsection{The result}

Let $p,d\ge1$, let $A_1,\dots,A_d$ be real $p\times p$ matrices, and let
\[
T_t=\sum_{k=1}^dA_kT_{t-k}+\xi_t\quad(t\ge1),\qquad T_t=0\quad(t\le0),
\]
with $\xi_t$ independent $\mathcal N(0,\Sigma_\xi)$, $\Sigma_\xi\succ0$.
Write $\Phi(z)=I-\sum_kA_kz^k$ and assume that $\det\Phi(z)\ne0$ on the
unit circle. The holder is that of Section~\ref{sec:model}. Let
$\varsigma=\max\{d,m+1\}$, set $A_k=0$ for $d<k\le\varsigma$, and let $\bA$
be the companion matrix of size $p\varsigma$ with first block row
$(A_1,\dots,A_\varsigma)$ and the identity below it. The state
$x_t=(T_t,\dots,T_{t-\varsigma+1})$ satisfies $x_{t+1}=\bA x_t+\bE\xi_{t+1}$
with $\bE=(I,0,\dots,0)^{\mathsf T}$, and $S_t=\bC x_t+U_t$ with
$\bC=(H_0,\dots,H_m,0,\dots,0)$. Since
$\det(\lambda I-\bA)=\lambda^{p\varsigma}\det\Phi(1/\lambda)$, the nonzero
eigenvalues of $\bA$ are the reciprocals of the zeros of $\det\Phi$; we
call them the poles. Let
\[
\mathcal N=\bigcap_{k\ge0}\ker(\bC\bA^k)
\]
be the unobservable subspace of the pair $(\bA,\bC)$, which is
invariant under $\bA$, and let
\[
\rho_\star=\sum_{|\rho|>1}\log|\rho|,\qquad
\rhid=\sum_{|\rho|>1,\ \rho\in\operatorname{spec}(\bA|_{\mathcal N})}\log|\rho|,
\]
the sums over eigenvalues with multiplicity, the first of $\bA$ and the
second of its restriction to $\mathcal N$. A pole of multiplicity one
has the eigenvector $(v,v/\rho,\dots,v/\rho^{\varsigma-1})$, up to a
factor, with $\Phi(1/\rho)v=0$, and $\bC$ maps it to $G(1/\rho)v$. By the test of
Popov, Belevitch, and Hautus such a pole is hidden, that is, an
eigenvalue of $\bA|_{\mathcal N}$, exactly when $G(1/\rho)v=0$: when the
holder's filter cancels it. Finally let
\[
f(\omega)=\Phi(e^{i\omega})^{-1}\Sigma_\xi\Phi(e^{i\omega})^{-*},
\]
which is continuous and positive definite, and let $V_f$ and
$\mathcal R_f(D)$ be as in Section~\ref{sec:model}, with
$\mathcal R_f(D)$ the nonnegative quadrant for $D\ge D_{\max}^f$.

\begin{theorem}\label{thm:vector}
For every $D>0$,
\[
\mathcal A(D)=\bigl(\rho_\star,\ \rhid\bigr)+\mathcal R_f(D).
\]
Deterministic encoders suffice, the converse holds with private
randomness, and for every $\alpha$ the window cost
$\tfrac1N\VG(\alpha,ND;\Gamma_N,\mathsf J_N)$ of $T^N$ and its holder
tends to $\alpha\rho_\star+(1-\alpha)\rhid+V_f(\alpha,D)$ for
$0<D<D_{\max}^f$, where $\mathsf J_N$ is the information matrix of the
holder's observations of the window, defined before
Lemma~\ref{lem:endpoint}.
\end{theorem}

For a scalar source and a memoryless holder with $H_0\ne0$, every pole
is observable: $\mathcal N$ contains at most eigenvectors of the
eigenvalue $0$ of the companion form, which exist when $A_d=0$. When
$H_0=0$, $\mathcal N=\R^{p\varsigma}$. So
$\rhid=\rho_\star\mathbf 1_{\{H_0=0\}}$ and Theorem~\ref{thm:vector} is
\cite[Theorem~26]{part3}. With the holder absent, the rate--distortion
function of the unstable vector source is $\rho_\star$ plus the
rate--distortion function of the stationary source with density $f$,
the vector form of \cite{hashimoto1980,grayhashimoto2008}. At
distortions small enough that reverse water-filling leaves every
frequency and every direction active, it equals
$\tfrac12\log\det(Q\Sigma_\xi)-\tfrac p2\log(D/p)$, because
$\int\log\det f\,\dw=\log\det\Sigma_\xi-2\rho_\star$. This is the
rate--distortion function of a single innovation vector, as Toms and
Berger found for the systems they treat \cite{toms1971}.

Instability is free in leakage exactly when $\rhid=0$, that is, when
every unstable mode is observable, which is the condition that the pair
$(\bA,\bC)$ be detectable. Three examples show what decides it. Let
$p=2$, $d=1$, $A_1=\operatorname{diag}(1.5,0.5)$, $\Sigma_\xi=I$, and a
memoryless holder. With $H_0=(1,0)$ the holder sees the unstable
coordinate, and the shift is $(\log1.5,0)$. With $H_0=(0,1)$ it sees
only the stable one, and the shift is $(\log1.5,\log1.5)$. With
$H_0=(\epsilon,1)$ the shift is $(\log1.5,0)$ for every $\epsilon\ne0$,
so $\rhid$ is not continuous in the holder: a mode seen with gain
$\epsilon$ grows like $1.5^t$ and rises above the noise after a time of
order $\log(1/\epsilon)$, which is negligible per symbol. Finally, for
the scalar source $T_t=\rho T_{t-1}+\xi_t$ with $|\rho|>1$, the holder
$S_t=T_t-\rho T_{t-1}+U_t=\xi_t+U_t$ observes the innovations. Its
filter $G(z)=1-\rho z$ vanishes at $1/\rho$, the pole is hidden, and
the shift is $(\log|\rho|,\log|\rho|)$: a holder that sees every
innovation, with noise, learns nothing about the growth.

\subsection{The window as a bridge and its final state}

Let $\Gamma_N=\Cov(T^N)$, let $\mathsf A_N$ be the block lower
triangular Toeplitz matrix with diagonal blocks $I$ and $k$th
subdiagonal blocks $-A_k$, so that $\mathsf A_NT^N=\xi^N$ and
$\Gamma_N=\mathsf A_N^{-1}(I_N\otimes\Sigma_\xi)\mathsf A_N^{-\mathsf T}$,
and let $\mathsf C_N$ be the full convolution matrix of $\Phi$, with
$N+d$ block rows: its first $N$ block rows are $\mathsf A_N$ and its
last $d$ block rows $\mathsf R_N$ give the terms
$(\Phi\ast T)_{N+j}$, $j=1,\dots,d$, which involve only
$T_{N-d+1},\dots,T_N$. For a trigonometric polynomial $g$ with
$p\times p$ coefficients, $T_N(g)$ is the block Toeplitz matrix whose
$(s,t)$ block is the coefficient of $e^{i(s-t)\omega}$ in $g$.

\begin{lemma}\label{lem:bridge}
Let $P_N=\mathsf C_N^{\mathsf T}(I_{N+d}\otimes\Sigma_\xi^{-1})\mathsf C_N$
and $B_N=\mathsf R_N^{\mathsf T}(I_d\otimes\Sigma_\xi^{-1})\mathsf R_N$.
Then $\Gamma_N^{-1}=P_N-B_N$, $B_N\succeq0$ has rank at most $pd$ and
is supported on the last $d$ blocks, and $P_N=T_N(f^{-1})$ with
$f^{-1}=\Phi(e^{i\omega})^*\Sigma_\xi^{-1}\Phi(e^{i\omega})$.
Let $X\sim\mathcal N(0,P_N^{-1})$, and split $T^N=(Y,Z)$ and $X=(Y_X,Z_X)$
into the first $N-d$ symbols and the last $d$. The conditional law of
$Y$ given $Z=z$ equals that of $Y_X$ given $Z_X=z$ for every $z$. It is
Gaussian with mean $m_Nz$, $m_N=-(P_N)_{YY}^{-1}(P_N)_{YZ}$, and
covariance $K_N=(P_N)_{YY}^{-1}=T_{N-d}(f^{-1})^{-1}$.
\end{lemma}

\begin{proof}
$\Gamma_N^{-1}=\mathsf A_N^{\mathsf T}(I_N\otimes\Sigma_\xi^{-1})\mathsf A_N=P_N-B_N$
by splitting the rows of $\mathsf C_N$. The $(s,t)$ block of $P_N$ is
$\sum_u\Phi_{u-s}^{\mathsf T}\Sigma_\xi^{-1}\Phi_{u-t}$, with
$\Phi_0=I$ and $\Phi_k=-A_k$, summed over the $u$ for which both factors
are nonzero; all such $u$ lie in $1,\dots,N+d$, so the block depends on
$s-t$ only and is the coefficient of $f^{-1}$ of that order.
Since $B_N$ vanishes outside the block of $Z$, the precision matrices
$\Gamma_N^{-1}$ and $P_N$ have the same blocks $YY$ and $YZ$, and the
conditional law of the first part of a Gaussian vector given the second
depends only on those blocks. The block $YY$ of $P_N$ is
$T_{N-d}(f^{-1})$.
\end{proof}

The density of $T^N$ times the density of the next $d$ innovations
evaluated at $T_{N+1}=\dots=T_{N+d}=0$ is proportional to
$\exp(-\tfrac12x^{\mathsf T}P_Nx)$, so $X$ is the source conditioned to
return to zero after the window: a Gaussian bridge. In the window the
source is this bridge together with its final state $Z$, and all of the
instability sits in the law of $Z$.

The bridge is close to a window of the stationary process with density
$f$. Since $f^{-1}$ is a positive definite trigonometric polynomial of
degree $d$, by the matrix form of the Fej\'er--Riesz theorem
\cite{rosenblum1985}, $f^{-1}=\tilde\Phi^*\tilde\Sigma^{-1}\tilde\Phi$ with
$\tilde\Phi(z)=I-\sum_{k\le d}\tilde A_kz^k$, $\det\tilde\Phi\ne0$ on the
closed unit disk, and $\tilde\Sigma\succ0$, and the stationary Gaussian
process $W$ with density $f$ is the causal autoregression with these
coefficients. On any window of times its density factors as the
density of its first $d$ values times the conditional densities of each
later value given the $d$ before. So its precision matrix on the window
is banded, and the block row of a time at distance at least $d$ from
both ends of the window collects all $d+1$ innovation terms and equals
the corresponding block row of the block Toeplitz matrix of $f^{-1}$.

\begin{lemma}\label{lem:stationary}
$K_N=\Cov(W_1,\dots,W_{N-d}\mid W_{-d+1},\dots,W_0,W_{N-d+1},\dots,W_N)$.
Consequently, for a group $\mathcal G$ of $n_0$ consecutive times in
$1,\dots,N-d$, the principal submatrix $\mathsf S_{\mathcal G}$ of $K_N$
on those times satisfies
$T_{n_0}(f)-c_W\,n_0\,\varrho^{2\operatorname{dist}(\mathcal G)}I\preceq\mathsf S_{\mathcal G}\preceq T_{n_0}(f)$,
where $T_{n_0}(f)=\Cov(W_1,\dots,W_{n_0})$, $\operatorname{dist}(\mathcal G)$
is the distance from $\mathcal G$ to $\{0,N-d+1\}$, and $c_W$ and
$\varrho<1$ depend only on $f$, $p$, and $d$.
\end{lemma}

\begin{proof}
On the window $-d+1,\dots,N$ the precision matrix of $W$ restricted to
the times $1,\dots,N-d$, all at distance at least $d$ from the ends, is
$T_{N-d}(f^{-1})$, and the inverse of a block of a precision matrix is
the conditional covariance given the rest. For the group,
$T_{n_0}(f)-\mathsf S_{\mathcal G}=\Cov(\E[W_{\mathcal G}\mid W_\partial])$,
with $W_\partial$ the $2d$ boundary values, which is positive
semidefinite with norm at most
$\|\Cov(W_{\mathcal G},W_\partial)\|^2\,\|\Cov(W_\partial)^{-1}\|$. The
covariances of a stationary process with a rational spectral density
and no poles on the circle decay geometrically,
$\|\Cov(W_s,W_t)\|\le c_0\varrho^{|s-t|}$, and
$\Cov(W_\partial)\succeq(\min_\omega\lambda_{\min}f(\omega))I$.
\end{proof}

\subsection{The holder's view of the growth}

\begin{lemma}\label{lem:innovations}
Let $R_t=\Cov(S_t\mid S^{t-1})$ be the innovation covariance of the
holder's observations. Then $R_t$ converges to a matrix $R_\infty$ with
\[
\log\det R_\infty=\int\log\det\bigl(\Sigma_U+G(e^{i\omega})f(\omega)G(e^{i\omega})^*\bigr)\dw+2(\rho_\star-\rhid).
\]
\end{lemma}

\begin{proof}
Choose coordinates in which the last ones span $\mathcal N$. Since
$\mathcal N$ is invariant and lies in $\ker\bC$,
\[
\bA=\begin{pmatrix}A_o&0\\A_{ho}&A_h\end{pmatrix},\qquad
\bC=(C_o,0),\qquad\bE=\begin{pmatrix}E_o\\E_h\end{pmatrix},
\]
with $(A_o,C_o)$ observable, and the spectrum of $\bA$ is that of $A_o$
together with that of $A_h=\bA|_{\mathcal N}$. The pair $(\bA,\bE)$ is
controllable, because the blocks of
$(\bE,\bA\bE,\dots,\bA^{\varsigma-1}\bE)$ form a block triangular matrix
with identity blocks on the diagonal, and the quotient pair
$(A_o,E_o)$ is controllable too. The observations
$S_t=C_ox_{o,t}+U_t$ depend only on the observable part, whose
dynamics $x_{o,t+1}=A_ox_{o,t}+E_o\xi_{t+1}$ do not involve the rest. So
$R_t=C_oP_tC_o^{\mathsf T}+\Sigma_U$, where $P_t=\Cov(x_{o,t}\mid S^{t-1})$
follows the Riccati recursion of the Kalman predictor of the observable
subsystem,
\[
P_{t+1}=A_oP_tA_o^{\mathsf T}+\bar Q-A_oP_tC_o^{\mathsf T}R_t^{-1}C_oP_tA_o^{\mathsf T},\qquad
\bar Q=E_o\Sigma_\xi E_o^{\mathsf T}.
\]

\emph{Convergence.} For $(A_o,C_o)$ detectable and $(A_o,\bar Q^{1/2})$
stabilizable, which hold here because the pairs are observable and
controllable, the recursion converges from every initial
$P_1\succeq0$ to the unique stabilizing solution $\bar P\succeq0$ of the
algebraic Riccati equation
$\bar P=A_o\bar PA_o^{\mathsf T}+\bar Q-KR_\infty K^{\mathsf T}$, with
$R_\infty=C_o\bar PC_o^{\mathsf T}+\Sigma_U$ and
$K=A_o\bar PC_o^{\mathsf T}R_\infty^{-1}$, and $A_o-KC_o$ has all
eigenvalues inside the unit circle
\cite[Ch.~4]{andersonmoore1979}, \cite{kailath2000}. This does not need
$A_o$ to be stable. Hence $R_t\to R_\infty$.

\emph{Factorization.} Let $\Psi(z)=(zI-A_o)^{-1}$ and
$\mathsf W(z)=I+C_o\Psi(z)K$. For $|z|=1$, $\bar z=1/z$, and
\[
\bar P-A_o\bar PA_o^{\mathsf T}=(zI-A_o)\bar P(\bar zI-A_o)^{\mathsf T}+(zI-A_o)\bar PA_o^{\mathsf T}+A_o\bar P(\bar zI-A_o)^{\mathsf T},
\]
as one sees by expanding the right side. Multiply the Riccati equation,
written as $\bar Q=\bar P-A_o\bar PA_o^{\mathsf T}+KR_\infty K^{\mathsf T}$,
on the left by $C_o\Psi(z)$ and on the right by its conjugate transpose,
and use $A_o\bar PC_o^{\mathsf T}=KR_\infty$:
\[
\Sigma_U+C_o\Psi\bar Q\Psi^*C_o^{\mathsf T}
=\Sigma_U+C_o\bar PC_o^{\mathsf T}+C_o\Psi KR_\infty+R_\infty K^{\mathsf T}\Psi^*C_o^{\mathsf T}+C_o\Psi KR_\infty K^{\mathsf T}\Psi^*C_o^{\mathsf T}
=\mathsf W R_\infty\mathsf W^*.
\]
This is an identity of rational functions on the circle, and it holds
whether or not $A_o$ is stable; when $A_o$ is stable its left side is
the spectral density of the holder's observation, and otherwise only
the identity is used. The transfer function from $\xi$ to
$\sum_jH_jT_{t-j}$ is $\bC(zI-\bA)^{-1}\bE z=C_o(zI-A_o)^{-1}E_oz$,
because the modes in $\mathcal N$ do not reach the output, so the left
side is $\Sigma_U+G(e^{i\omega})f(\omega)G(e^{i\omega})^*$ up to the
conjugation $\omega\mapsto-\omega$, which does not change its integral.

\emph{Determinant and Jensen.} By the identity
$\det(I+XY)=\det(I+YX)$,
$\det\mathsf W(z)=\det(I+\Psi(z)KC_o)=\det(zI-A_o+KC_o)/\det(zI-A_o)$.
For a square matrix $M$ with eigenvalues $\lambda_i$,
$\det(e^{i\omega}I-M)=\prod_i(e^{i\omega}-\lambda_i)$, and
$\int\log|e^{i\omega}-\lambda|\dw=\log\max\{1,|\lambda|\}$ by Jensen's
formula, so
$\int\log|\det(e^{i\omega}I-M)|\dw=\sum_{|\lambda_i|>1}\log|\lambda_i|$.
Neither $A_o$ nor $A_o-KC_o$ has eigenvalues on the circle, the latter
has none outside it, and the eigenvalues of $A_o$ outside the circle
are those of $\bA$ outside it that do not belong to $\bA|_{\mathcal N}$.
Hence $\int\log|\det\mathsf W|\dw=-(\rho_\star-\rhid)$. Taking $\log\det$
of the factorization and integrating gives the claim.
\end{proof}

The next lemma uses block circulant matrices. For a trigonometric
polynomial $g$ of degree $L<N/2$, $C_N(g)$ is the block circulant matrix
whose $(s,t)$ block is the sum of the coefficients of $g$ of the orders
congruent to $s-t$ modulo $N$. It differs from $T_N(g)$ only in the
blocks with $|s-t|\ge N-L$, a matrix of rank at most $2pL$ whose norm is
bounded by the sum of the norms of the coefficients, and the discrete
Fourier transform makes it block diagonal with the blocks $g$ at the
$N$ Fourier frequencies. In particular
$\log\det C_N(g)=N\int\log\det g\,\dw+o(N)$ for a positive definite $g$,
a Riemann sum of a continuous function. Let $\mathsf G_N$ be the matrix
of the holder on the window, with the blocks $H_{s-t}$ for
$0\le s-t\le m$, and
$\mathsf J_N=\mathsf G_N^{\mathsf T}(I_N\otimes\Sigma_U^{-1})\mathsf G_N$.
As for $P_N$, $\mathsf J_N=T_N(j)$ minus a positive semidefinite matrix
of rank at most $rm$ on the last $m$ blocks, so
$\mathsf J_N-C_N(j)$ has rank at most $rm+2pm$ and bounded norm.

\begin{lemma}\label{lem:endpoint}
\begin{enumerate}
\item[(a)] $\log\det\Cov(Z)=2N\rho_\star+o(N)$.
\item[(b)] $\log\det\Cov(Z\mid S^N)=2N\rhid+o(N)$.
\item[(c)] Let $m_f>0$ and $M_f$ be the smallest and the largest
eigenvalue of $f(\omega)$ over $\omega$,
$\bar J=(\sum_j\|H_j\|)^2\|\Sigma_U^{-1}\|$, and
$c=\min\{(1/m_f+\bar J)^{-1},1/M_f\}$, which does not depend on $N$.
With $S_X=\mathsf G_NX+U$, the four matrices $\Cov(Z_X)$,
$\Cov(Z_X\mid S_X)$, $\Cov(Z)$, and $\Cov(Z\mid S^N)$ are all at least
$cI$, and the first two are at most $c^{-1}I$.
\end{enumerate}
By (a) and (b) the volumes of the last two grow at the rates $\rho_\star$
and $\rhid$, so they have no upper bound independent of $N$ when these
rates are positive.
\end{lemma}

\begin{proof}
(c) For a block Toeplitz matrix with a Hermitian symbol $g$,
$x^*T_N(g)x=\int\hat x(\omega)^*g(\omega)\hat x(\omega)\dw$ with
$\int\|\hat x\|^2\dw=\|x\|^2$, so its eigenvalues lie between the extreme
eigenvalues of $g$. Hence $P_N=T_N(f^{-1})$ has eigenvalues in
$[1/M_f,1/m_f]$ and $P_N^{-1}$ in $[m_f,M_f]$. The matrix $\Cov(Z_X)$ is
a principal submatrix of $P_N^{-1}$, so its eigenvalues lie in
$[m_f,M_f]$ by interlacing, and $\Cov(Z_X\mid S_X)\preceq\Cov(Z_X)$.
For the lower bounds, conditioning on more lowers the covariance.
Given $Y$ and $S^N$, the conditional precision of $Z$ is the block $ZZ$
of the posterior precision $\Gamma_N^{-1}+\mathsf J_N$ of $T^N$, so
\[
\Cov(Z)\succeq\Cov(Z\mid S^N)\succeq\Cov(Z\mid Y,S^N)=\bigl((\Gamma_N^{-1})_{ZZ}+(\mathsf J_N)_{ZZ}\bigr)^{-1}.
\]
By Lemma~\ref{lem:bridge}, $(\Gamma_N^{-1})_{ZZ}=(P_N)_{ZZ}-B_N$, and
$B_N\succeq0$, so $(\Gamma_N^{-1})_{ZZ}\preceq(P_N)_{ZZ}\preceq(1/m_f)I$;
and $\|\mathsf J_N\|\le\|\mathsf G_N\|^2\|\Sigma_U^{-1}\|\le\bar J$. So
the right side is at least $(1/m_f+\bar J)^{-1}I$. This bound does not
depend on $N$, however fast the hidden modes grow: growth makes
$(\Gamma_N^{-1})_{ZZ}$ smaller, never larger. The same argument with
$P_N$ in place of $\Gamma_N^{-1}$ gives the bounds for $Z_X$.

(a) By Lemma~\ref{lem:bridge} the conditional laws given the final
blocks agree, so
$\log\det\Gamma_N-\log\det\Cov(Z)=\log\det K_N=\log\det P_N^{-1}-\log\det\Cov(Z_X)$.
Since $\det\mathsf A_N=1$, $\log\det\Gamma_N=N\log\det\Sigma_\xi$. The
matrices $P_N=T_N(f^{-1})$ and $C_N(f^{-1})$ have eigenvalues in
$[1/M_f,1/m_f]$ and differ by a matrix of rank at most $2pd$ and bounded
norm, so their log-determinants differ by $O(1)$, and
$\log\det C_N(f^{-1})=N\int\log\det f^{-1}\dw+o(N)$. Finally
$\int\log\det f\,\dw=\log\det\Sigma_\xi-2\int\log|\det\Phi(e^{i\omega})|\dw$,
and Jensen's formula for the polynomial $\det\Phi$, with
$\det\Phi(0)=1$, gives $\int\log|\det\Phi|\dw=\sum\log(1/|z_j|)$ over its
zeros in the unit disk, which is $\rho_\star$. With (c),
$\log\det\Cov(Z)=2N\rho_\star+o(N)$.

(b) In bits,
$\log\det\Cov(Z\mid S^N)=\log\det\Cov(Z)-2I(Z;S^N)$. Write
$I(T^N;S^N)=I(Z;S^N)+I(Y;S^N\mid Z)$ and the same for $X$. Given $Z=z$
the pair $(Y,S^N)$ is Gaussian with a mean that depends on $z$ and a
covariance that does not, and by Lemma~\ref{lem:bridge} that covariance
is the same for $X$ given $Z_X=z$, so the two conditional informations
are equal. Hence
\[
I(Z;S^N)=I(Z_X;S_X)+I(T^N;S^N)-I(X;S_X).
\]
By (c), $I(Z_X;S_X)=O(1)$. Next,
$I(X;S_X)=\tfrac12[\log\det(P_N+\mathsf J_N)-\log\det P_N]$, and
$P_N+\mathsf J_N$ differs from $C_N(f^{-1}+j)$ by a matrix of rank at
most $2pd+rm+2pm$ and bounded norm, both having eigenvalues in a fixed
interval, so
$\tfrac1NI(X;S_X)\to\tfrac12\int\log\det(I+fj)\dw$. Finally
$I(T^N;S^N)=\sum_t[h(S_t\mid S^{t-1})-h(U_t)]=\tfrac12\sum_t\log\det(R_t\Sigma_U^{-1})$,
and by Lemma~\ref{lem:innovations} and the identity
$\det(I+\Sigma_U^{-1}GfG^*)=\det(I+fj)$,
\[
\tfrac1NI(T^N;S^N)\to\tfrac12\int\log\det(I+fj)\dw+\rho_\star-\rhid.
\]
Hence $I(Z;S^N)=N(\rho_\star-\rhid)+o(N)$, and (a) gives (b).
\end{proof}

\subsection{The window cost}

\begin{lemma}\label{lem:perturb}
For $\Pi\succ0$ and $\mathsf J\succeq0$ of size $pn$ let
$\mathsf c(\Pi,\mathsf J)=\tfrac\alpha2\log\det\Pi+\tfrac{1-\alpha}2\log\det(\Pi+\mathsf J)$.
Let $\Pi,\Pi'\succ0$ and $\mathsf J,\mathsf J'\succeq0$, with $\Pi'-\Pi$ and
$\mathsf J'-\mathsf J$ of rank at most $\bar k$ and norm at most $b$.
Then for every $D>0$
\[
\Bigl|\bigl[\VG(\alpha,nD;\Pi^{-1},\mathsf J)+\mathsf c(\Pi,\mathsf J)\bigr]-\bigl[\VG(\alpha,nD;\Pi'^{-1},\mathsf J')+\mathsf c(\Pi',\mathsf J')\bigr]\Bigr|
\le\bar k\log\Bigl(1+\frac{bnD}q\Bigr).
\]
\end{lemma}

\begin{proof}
Write $P=\Delta^{-1}$ and $\Lambda=\log(1+bnD/q)$. In \eqref{eq:costs},
$\alpha r+(1-\alpha)\ell=g_{\mathsf J}(P)-\mathsf c(\mathsf S^{-1},\mathsf J)$
with
$g_{\mathsf J}(P)=\tfrac\alpha2\log\det P+\tfrac{1-\alpha}2\log\det(P+\mathsf J)$,
and the constraints are $P\in\mathcal F(\mathsf S^{-1})$, where
$\mathcal F(\Pi)=\{P\succeq\Pi:\ \tr((I_n\otimes Q)P^{-1})\le nD\}$. So the
bracket for $(\Pi,\mathsf J)$ is $\min_{\mathcal F(\Pi)}g_{\mathsf J}$. On
$\mathcal F(\Pi)$, $q\,\lambda_{\max}(P^{-1})\le nD$, so
$\lambda_{\min}(P)\ge q/(nD)$, and for a positive semidefinite $B$ of
rank $k'$ and norm at most $b$,
$0\le\log\det(P+B)-\log\det P\le k'\Lambda$, and the same with
$P+\mathsf J$ in place of $P$. For a symmetric $X$ of rank at most
$\bar k$, its positive and negative parts $X_\pm\succeq0$ have norms at
most $\|X\|$ and ranks with $\operatorname{rank}X_++\operatorname{rank}X_-\le\bar k$.

\emph{The information matrix.} Let $\mathsf J'-\mathsf J=B_+-B_-$ be that
decomposition, so that $\mathsf J+B_+=\mathsf J'+B_-$. Then
$g_{\mathsf J}\le g_{\mathsf J+B_+}\le g_{\mathsf J}+\tfrac{\bar k}2\Lambda$ and
$g_{\mathsf J'}\le g_{\mathsf J'+B_-}\le g_{\mathsf J'}+\tfrac{\bar k}2\Lambda$
on every $\mathcal F(\Pi)$. Both $g_{\mathsf J}$ and $g_{\mathsf J'}$ lie
within $\tfrac{\bar k}2\Lambda$ below the same function, so
$|g_{\mathsf J}-g_{\mathsf J'}|\le\tfrac{\bar k}2\Lambda$.

\emph{The precision.} Let $\Pi''=\Pi+(\Pi'-\Pi)_+=\Pi'+(\Pi'-\Pi)_-$,
which dominates $\Pi$ and $\Pi'$, so
$\mathcal F(\Pi'')\subseteq\mathcal F(\Pi)\cap\mathcal F(\Pi')$ and
$\min_{\mathcal F(\Pi'')}g$ is at least each of $\min_{\mathcal F(\Pi)}g$
and $\min_{\mathcal F(\Pi')}g$. Conversely, for $P\in\mathcal F(\Pi)$ the
matrix $P+(\Pi'-\Pi)_+$ dominates $\Pi''$, its inverse is at most
$P^{-1}$, so it lies in $\mathcal F(\Pi'')$, and its cost is at most
$g(P)+\tfrac{\bar k}2\Lambda$; the same holds with $\Pi'$ and
$(\Pi'-\Pi)_-$. So $\min_{\mathcal F(\Pi)}g$ and $\min_{\mathcal F(\Pi')}g$
both lie within $\tfrac{\bar k}2\Lambda$ below $\min_{\mathcal F(\Pi'')}g$,
and they differ by at most $\tfrac{\bar k}2\Lambda$.

The two steps, applied one after the other, give the claim.
\end{proof}

\begin{lemma}\label{lem:circulant}
For $0<D<D_{\max}^f$,
$\tfrac1N\VG\bigl(\alpha,ND;C_N(f^{-1})^{-1},C_N(j)\bigr)\to V_f(\alpha,D)$.
\end{lemma}

\begin{proof}
The cyclic shift of the window by one symbol commutes with
$C_N(f^{-1})$, $C_N(j)$, and $I_N\otimes Q$, so it maps the feasible set
to itself and preserves the cost. The cost is convex by
\cite[Theorem~5]{part1}, so the average of a minimizer over the $N$
shifts is a block circulant minimizer. The discrete Fourier transform
makes every block circulant matrix block diagonal, and the problem
becomes the program of Section~\ref{sec:model} over the uniform measure
on the $N$ Fourier frequencies, with the covariances
$f(2\pi k/N)$ and the information matrices $j(2\pi k/N)$; a minimizer
may be taken conjugate symmetric across $k$ and $N-k$, hence real
after the transform. By (ii) of Section~\ref{sec:model} its cost tends
to $V_f(\alpha,D)$.
\end{proof}

\subsection{Codes for groups of symbols}

\begin{lemma}\label{lem:groups}
Let $n_0>m$, and let $\mathsf G'$ be the map from $n_0$ symbols to the
holder's observations at their times $m+1,\dots,n_0$, which involve only
those symbols, with information matrix $\mathsf J'$. Let a code with a
deterministic encoder for the Gaussian vector $T'$ of $n_0$ symbols with
covariance $\mathsf S'\succ0$ have rate $R'$, distortion $D'$, and
leakage $L'=\tfrac1{n_0}H(M'\mid\mathsf G'T'+U')$, with $U'$ white noise of
covariance $I\otimes\Sigma_U$ independent of $T'$. Let $T$ be a Gaussian
vector whose $k$ consecutive groups $T_1,\dots,T_k$ of $n_0$ symbols,
possibly dependent, have covariances $\mathsf S_j$ with
$\mathsf S'\preceq\mathsf S_j\preceq b_0I$. Let $S$ be any observation that
contains, for each $j$, the vector $S_j=\mathsf G'T_j+U_j$, where the $U_j$
are white noises of covariance $I\otimes\Sigma_U$, independent of each
other, of $T$, and of the rest of $S$ given $T$. Then for every
$\delta>0$ and every $k\ge k_0$, with $k_0$ depending only on the code,
$b_0$, $Q$, and $\delta$, there is a code for $T$ with a deterministic
encoder, rate $R'$, distortion at most $D'+\epsilon_Q+\delta$, and
$\tfrac1{kn_0}H(M\mid S)\le L'+\epsilon_J+\delta$, where
$\epsilon_Q=\max_j\tfrac1{n_0}\tr((I\otimes Q)(\mathsf S_j-\mathsf S'))$ and
$\epsilon_J=\max_j\tfrac1{2n_0\ln2}\tr(\mathsf J'(\mathsf S_j-\mathsf S'))$.
\end{lemma}

\begin{proof}
This is \cite[Lemma~12]{part3}, whose proof uses the holder in three
places. First, the coupling of \cite[Lemma~11]{part3} draws $T'_j$ from
$T_j$ so that $(T'_j,\mathsf G'T'_j+U_j)$ has the law under which $L'$ was
measured; this holds here because $S_j$ is $\mathsf G'T_j$ plus white
noise of the right covariance. Second, the leakage the coupling adds is
at most $I(E_j;\mathsf G'E_j+U_j)\le\tr(\mathsf J'\Cov E_j)/(2\ln2)$, which
holds for any linear map $\mathsf G'$. Third, the leakage of the tuple of
messages is bounded by
$H(M\mid S)\le\sum_jH(M_j\mid S)\le\sum_jH(M_j\mid S_j)$, which holds
because each $S_j$ is part of $S$.
\end{proof}

A code from (i) of Section~\ref{sec:model} for the stationary source
with density $f'$ has its leakage measured against the holder's
observations at all the times of its window, $A=(S_1,\dots,S_m)$ and
$B=(S_{m+1},\dots,S_{n_0})$, and $A$ involves up to $m$ symbols before the
window. Its leakage against $B$ alone is larger by
$\tfrac1{n_0}I(M;A\mid B)$. Since $M$ is a function of the window and
$A$ and $B$ are independent given the source on the times
$1-m,\dots,n_0$, $I(M;A\mid B)\le I(W_{1-m..n_0};A\mid B)\le I(W_{1-m..n_0};A)$,
the information through a channel of dimension $rm$ with white noise.
By Fischer's inequality applied to the covariance of $A$ it is at most
$\tfrac{rm}2\log(1+\kappa_0)$, with $\kappa_0$ the largest eigenvalue of
$\Sigma_U^{-1/2}\Cov(\sum_jH_jW_{t-j})\Sigma_U^{-1/2}$, which is at most
$\bar JM_f$. So against the observations of Lemma~\ref{lem:groups} the
code's leakage per symbol is larger by at most
$\tfrac{rm}{2n_0}\log(1+\bar JM_f)$, which tends to zero as $n_0$ grows.

\subsection{A finite quantizer for the final state}

The rate of a code is the logarithm of the size of its message set, so
the final state, whose variance grows exponentially, must be described
with finitely many messages.

\begin{lemma}\label{lem:quantizer}
Let $Z_N\in\R^k$ be centered and jointly Gaussian with $S^N$, with
covariance $\Xi_N$, and let $c,\beta>0$ be such that $\Xi_N\succeq cI$,
$\Cov(Z_N\mid S^N)\succeq cI$, and $\ln\lambda_{\max}(\Xi_N)\le\beta N$ for
every $N$. For every $\epsilon>0$ there are deterministic quantizers
$M_{Z,N}=\mathsf q_N(Z_N)$ with finite index sets $\mathcal M_{Z,N}$ and
reproductions $\hat Z_N$ such that
\begin{enumerate}
\item[(a)] $\log|\mathcal M_{Z,N}|\le\tfrac12\log\det\Xi_N+O(\log N)$;
\item[(b)] $H(M_{Z,N}\mid S^N)\le\tfrac12\log\det\Cov(Z_N\mid S^N)+O(1)$;
\item[(c)] $\E\|Z_N-\hat Z_N\|^2\le k\epsilon^2/4+o(1)$,
\end{enumerate}
with constants that depend on $k$, $c$, $\beta$, and $\epsilon$ only.
\end{lemma}

\begin{proof}
Let $u_1,\dots,u_k$ be orthonormal eigenvectors of $\Xi_N$ with
eigenvalues $\lambda_i$, let $V_i=u_i^{\mathsf T}Z_N$, and let
$s_N=(2\beta N+4\ln N)^{1/2}$. Let $j_i=\lfloor V_i/\epsilon\rfloor$ be
the index of the lattice cell of $V_i$ and $J_i=\lceil s_N\sqrt{\lambda_i}/\epsilon\rceil$.
The quantizer sends $j_i$ when $|j_i|\le J_i$ and an overflow symbol
otherwise, and the reproduction is the midpoint $(j_i+\tfrac12)\epsilon$
of the cell, or $0$ after an overflow, in each coordinate.

(a) The index set has $\prod_i(2J_i+2)$ elements, and
$2J_i+2\le2s_N\sqrt{\lambda_i}/\epsilon+4\le\sqrt{\lambda_i}\,(2s_N/\epsilon+4/\sqrt c)$
because $\lambda_i\ge c$. So
$\log|\mathcal M_{Z,N}|\le\tfrac12\log\det\Xi_N+k\log(2s_N/\epsilon+4/\sqrt c)$,
and $s_N=O(\sqrt N)$.

(b) The index is a function of the vector $\mathbf j=(j_1,\dots,j_k)$ of
lattice indices, so $H(M_{Z,N}\mid S^N)\le H(\mathbf j\mid S^N)$. Let
$\mathsf w$ be uniform on $[0,\epsilon)^k$ and independent of
$(Z_N,S^N)$. Given $S^N$, the vector $\epsilon\mathbf j+\mathsf w$ has the
conditional density of $\mathbf j$ spread evenly over the cells, so
$h(\epsilon\mathbf j+\mathsf w\mid S^N)=H(\mathbf j\mid S^N)+k\log\epsilon$.
It equals $V+\Delta$ with $V=(V_i)$ and $\Delta\in(-\epsilon,\epsilon)^k$,
so its conditional covariance is at most $2\Cov(V\mid S^N)+2k\epsilon^2I$.
By the maximum of entropy at a given covariance and the concavity of
$\log\det$,
\begin{equation}
H(\mathbf j\mid S^N)\le\tfrac12\log\det\bigl(2\pi e\cdot2(\Cov(Z_N\mid S^N)+k\epsilon^2I)\bigr)-k\log\epsilon,
\label{eq:quant}
\end{equation}
since $\Cov(V\mid S^N)$ is $\Cov(Z_N\mid S^N)$ in the basis $(u_i)$. With
$\Cov(Z_N\mid S^N)\succeq cI$ the right side is
$\tfrac12\log\det\Cov(Z_N\mid S^N)+\tfrac k2\log(4\pi e(1+k\epsilon^2/c))-k\log\epsilon$.

(c) Without overflow the error in each coordinate is at most
$\epsilon/2$. Overflow in coordinate $i$ requires
$|V_i|\ge s_N\sqrt{\lambda_i}$, and for a standard Gaussian variable
$\mathsf g$,
$\E[\mathsf g^2\mathbf 1_{\{|\mathsf g|\ge s\}}]=2s\varphi(s)+2\bar\Phi(s)\le(s+1)e^{-s^2/2}$.
So the error from overflow is at most
$\sum_i\lambda_i(s_N+1)e^{-s_N^2/2}\le k(s_N+1)/N^2$, because
$\lambda_i\le e^{\beta N}$.
\end{proof}

For the source of Theorem~\ref{thm:vector}, with $Z=Z_N$ the final state
and $k=pd$, the hypotheses hold: the two lower bounds by
Lemma~\ref{lem:endpoint}(c), and the upper one because
$T_t=\sum_{j<t}\Psi_j\xi_{t-j}$ with impulse response
$\|\Psi_j\|\le c'(1+j)^{p\varsigma}\bar\rho^{\,j}$, where $\bar\rho$ is the
largest modulus of a pole, so that $\ln\lambda_{\max}(\Xi_N)\le\beta N$ for
every $N\ge1$ with $\beta$ large enough.

\subsection{Proof of Theorem~\ref{thm:vector}}

\emph{The window cost.} Let $D>0$. Apply Lemma~\ref{lem:perturb} with
$\Pi=\Gamma_N^{-1}=P_N-B_N$, $\mathsf J=\mathsf J_N$, $\Pi'=C_N(f^{-1})$,
and $\mathsf J'=C_N(j)$. The differences have ranks at most $3pd$ and
$rm+2pm$, so $\bar k=\max\{3pd,rm+2pm\}$, and norms bounded in $N$, so
\[
\VG(\alpha,ND;\Gamma_N,\mathsf J_N)=\VG\bigl(\alpha,ND;C_N(f^{-1})^{-1},C_N(j)\bigr)+\mathsf c\bigl(C_N(f^{-1}),C_N(j)\bigr)-\mathsf c(\Gamma_N^{-1},\mathsf J_N)+O(\log N)
\]
for every $D>0$. For the constants, $\log\det\Gamma_N^{-1}=-N\log\det\Sigma_\xi$
and
$\log\det C_N(f^{-1})=N\int\log\det f^{-1}\dw+o(N)=N(2\rho_\star-\log\det\Sigma_\xi)+o(N)$,
by the computation in the proof of Lemma~\ref{lem:endpoint}(a), so the
first difference of log-determinants is $2N\rho_\star+o(N)$. For the
second,
$\log\det(\Gamma_N^{-1}+\mathsf J_N)=\log\det\Gamma_N^{-1}+2I(T^N;S^N)$,
which by the proof of Lemma~\ref{lem:endpoint}(b) is
$N[-\log\det\Sigma_\xi+\int\log\det(I+fj)\dw+2(\rho_\star-\rhid)]+o(N)$,
while
$\log\det(C_N(f^{-1})+C_N(j))=N\int\log\det(f^{-1}+j)\dw+o(N)=N[2\rho_\star-\log\det\Sigma_\xi+\int\log\det(I+fj)\dw]+o(N)$.
The second difference is therefore $2N\rhid+o(N)$, and
\[
\mathsf c\bigl(C_N(f^{-1}),C_N(j)\bigr)-\mathsf c(\Gamma_N^{-1},\mathsf J_N)=N[\alpha\rho_\star+(1-\alpha)\rhid]+o(N).
\]
For $0<D<D_{\max}^f$, Lemma~\ref{lem:circulant} gives the limit of the
window cost. For every $D>0$ the display and $\VG\ge0$ give
$\liminf\tfrac1N\VG(\alpha,ND;\Gamma_N,\mathsf J_N)\ge\alpha\rho_\star+(1-\alpha)\rhid$.

\emph{Converse.} Let $(R,L)$ be achievable at $D$ with codes of
blocklengths $N_j$, and let $D'>D$. Eventually $D_{N_j}\le D'$, so by
Section~\ref{sec:model}
$\alpha R_{N_j}+(1-\alpha)L_{N_j}\ge\tfrac1{N_j}\VG(\alpha,N_jD';\Gamma_{N_j},\mathsf J_{N_j})$.
In the limit,
$\alpha R+(1-\alpha)L\ge\alpha\rho_\star+(1-\alpha)\rhid+V_f(\alpha,D')$,
by the limit of the window cost for $D'<D_{\max}^f$ and by its lower
bound, with $V_f=0$, for $D'\ge D_{\max}^f$. Letting $D'\downarrow D$ and
using the continuity of $V_f$, $(R,L)$ lies in every half-plane that
defines the translate of $\mathcal R_f(D)$.

\emph{Achievability: the code.} Let $0<\eta<m_f$, $f'=f-\eta I$, and let
$(R',L')\in\mathcal R_{f'}(D')$ with $0<D'<D_{\max}^{f'}=D_{\max}^f-\eta\tr Q$.
By (i) of Section~\ref{sec:model} and the remark after
Lemma~\ref{lem:groups}, there is, for every $\delta>0$ and some large
$n_0>m$, a code of blocklength $n_0$ for the stationary source with
density $f'$ and the holder, with a deterministic encoder, rate at most
$R'+\delta$, distortion at most $D'+\delta$, and leakage at most
$L'+\delta$ against the observations at its times after the first $m$.
For a blocklength $N$ the encoder does three things.
\begin{enumerate}
\item It quantizes $Z$ with the quantizer of Lemma~\ref{lem:quantizer},
for a fixed $\epsilon>0$, and sends the index $M_Z$; the decoder takes
its reproduction $\hat Z$.
\item It computes $Y''=Y-m_NZ$, which by Lemma~\ref{lem:bridge} is
$\mathcal N(0,K_N)$ and independent of $Z$. Let $L_\eta$, which depends
on $\eta$ and $n_0$, be such that
$c_W\,n_0\,\varrho^{2L_\eta}\le\eta$ in Lemma~\ref{lem:stationary}. The
times of $Y''$ at distance at least $L_\eta$ from both ends are cut into
$k_N$ groups of $n_0$, with $k_N\to\infty$; by Lemma~\ref{lem:stationary} each group has
covariance $\mathsf S_{\mathcal G}$ with
$T_{n_0}(f')\preceq\mathsf S_{\mathcal G}\preceq T_{n_0}(f)$, and the code
for $f'$ is applied to the groups by Lemma~\ref{lem:groups}, which adds
at most $\eta\tr Q$ in distortion and $\eta\bar Jp/(2\ln2)$ in leakage per
symbol, plus $\delta$. The other times of $Y''$, at most $2L_\eta+n_0$ of
them, are reproduced as zero.
\item The decoder reproduces $\hat Y=\hat Y''+m_N\hat Z$.
\end{enumerate}
The encoder is deterministic and the message set is the product of
$\mathcal M_{Z,N}$ and the message sets of the group codes, all finite.

\emph{Distortion.} The error on $Y$ is $(Y''-\hat Y'')+m_N(Z-\hat Z)$,
and $\|m_N\|$ is bounded in $N$ because $P_N$ has eigenvalues in a
fixed interval. So for any $\zeta>0$ the distortion per symbol is at
most
$(1+\zeta)\tfrac1N\E\|Y''-\hat Y''\|_Q^2+(2+1/\zeta)\|Q\|(1+\|m_N\|^2)\,\E\|Z-\hat Z\|^2/N$.
The second term is $O(1/N)$ by Lemma~\ref{lem:quantizer}(c), and the
first tends to at most $(1+\zeta)(D'+2\delta+\eta\tr Q)$, the share of
the zero-coded times tending to zero.

\emph{Rate.} By Lemma~\ref{lem:quantizer}(a) and
Lemma~\ref{lem:endpoint}(a), $\tfrac1N\log|\mathcal M_{Z,N}|\to\rho_\star$,
and the groups add at most $R'+\delta$ per symbol.

\emph{Leakage.} With $M=(M_Z,M_Y)$,
$H(M\mid S^N)\le H(M_Z\mid S^N)+H(M_Y\mid S^N)$. By
Lemma~\ref{lem:quantizer}(b) and Lemma~\ref{lem:endpoint}(b),
$\limsup\tfrac1NH(M_Z\mid S^N)\le\rhid$. For the other term, whiten the
noise, multiplying the observations by $I_N\otimes\Sigma_U^{-1/2}$, and
write the whitened observations as $\tilde S''+\tilde{\mathsf G}_ZZ$,
where $\tilde S''$ is the whitened $\mathsf G_N(Y'',0)+U^N$ and
$\tilde{\mathsf G}_Z$, of rank at most $pd$, gives the contribution of
$Z$ through $m_NZ$ and through $Z$ itself. Let $\Pi_Z$ be the orthogonal
projection onto the complement of the range of $\tilde{\mathsf G}_Z$, and
$A=(I-\Pi_Z)\tilde S''$, $B=\Pi_Z\tilde S''$. The whitened noise is white,
so its two projections are independent, and $A$ and $B$ are independent
given $Y''$. Since $M_Y$ is a function of $Y''$ and $Y''$, $Z$, and the
noise are independent,
\[
H(M_Y\mid S^N)\le H(M_Y\mid B)=H(M_Y\mid A,B)+I(M_Y;A\mid B)\le H(M_Y\mid\tilde S'')+I(Y'';A),
\]
using $I(M_Y;A\mid B)\le I(Y'';A\mid B)\le I(Y'';A)$, the last because
$A$ and $B$ are independent given $Y''$. The term $I(Y'';A)$ is the
information through a channel of dimension at most $pd$ with bounded
signal and white noise, which is $O(1)$. Finally
Lemma~\ref{lem:groups} bounds $H(M_Y\mid\tilde S'')$, which equals the
same quantity without whitening. So the leakage per symbol is at most
$\rhid+L'+\eta\bar Jp/(2\ln2)+2\delta+o(1)$.

\emph{Closing the argument.} Let $0<D<D_{\max}^f$, $(R,L)\in\mathcal R_f(D)$,
and $\gamma>0$. The function $V_f(\alpha,\cdot)$ is convex and
nonincreasing, and $V_f(\alpha,D/2)\le V_f(1,D/2)$ because $\ell\le r$,
so its slopes on $[D/2,D]$ are bounded uniformly in $\alpha$; choose
$D''\in[D/2,D)$ with $V_f(\alpha,D'')\le V_f(\alpha,D)+\gamma$ for every
$\alpha$. Then choose $\eta>0$ so small that, with
$\kappa_\eta=\eta/(m_f-\eta)$, $\tfrac p2\log(1+\kappa_\eta)\le\gamma$, and
$\zeta,\delta>0$ so small that $(1+\zeta)(D'+2\delta+\eta\tr Q)<D$ for
$D'=D''/(1+\kappa_\eta)$; this is possible because $D'<D''<D$. Since
$f'\preceq f\preceq(1+\kappa_\eta)f'$, the left inequality of (iii),
evaluated at $(1+\kappa_\eta)D'=D''$, gives
$V_{f'}(\alpha,D')\le V_f(\alpha,D'')+\gamma\le V_f(\alpha,D)+2\gamma$ for
every $\alpha$, and $D'<D<D_{\max}^f-\eta\tr Q$ once $\eta$ is small. As
$\alpha R+(1-\alpha)L\ge V_f(\alpha,D)$ for every $\alpha$, the pair
$(R+2\gamma,L+2\gamma)$ lies in every half-plane that defines
$\mathcal R_{f'}(D')$, hence in it. The code above with
$(R',L')=(R+2\gamma,L+2\gamma)$ has, for large $N$, distortion below $D$,
rate at most $\rho_\star+R+2\gamma+2\delta$, and leakage at most
$\rhid+L+2\gamma+\eta\bar Jp/(2\ln2)+3\delta$. Letting $\gamma$, $\eta$,
and $\delta$ tend to zero puts $(\rho_\star+R,\rhid+L)$ in the closure
$\mathcal A(D)$. For $D\ge D_{\max}^f$ only the index of
Lemma~\ref{lem:quantizer} is sent and $Y$ is reproduced as $m_N\hat Z$,
whose distortion tends to at most $D_{\max}^f$ because
$K_N\preceq T_{N-d}(f)$, which gives $(\rho_\star,\rhid)$ and every pair
above it.
\hfill$\blacksquare$

\section{Sources with a Limiting Eigenvalue Distribution}
\label{sec:law}

\subsection{The result}

Let $(T_t)_{t\ge1}$ be a centered scalar Gaussian process, with any
covariance, observed by the holder through $S_t=HT_t+U_t$, $H\in\R$, with
$U_t$ independent $\mathcal N(0,\sigma_U^2)$, and $J=H^2/\sigma_U^2$. The
decoder reproduces $T$ with weight $Q>0$. Let $\Gamma_N=\Cov(T^N)$ and
let $\mu_N$ be the distribution of its eigenvalues, the probability
measure that puts mass $1/N$ on each.

\begin{theorem}\label{thm:law}
Assume that $\|\Gamma_N\|\le b$ for every $N$ and that $\mu_N$ converges
weakly to a probability measure $\mu$. Let $D_{\max}^\mu=Q\int\lambda\,d\mu$,
and for $0<D<D_{\max}^\mu$ let $\theta>0$ solve
$Q\int\min\{\lambda,\theta\}\,d\mu=D$ and
\[
r_\mu(D)=\int\tfrac12\log^+\frac\lambda\theta\,d\mu,\qquad
\ell_\mu(D)=\int_{\lambda>\theta}\tfrac12\log\frac{1/\theta+J}{1/\lambda+J}\,d\mu .
\]
Then $\mathcal A(D)$ is the quadrant above $(r_\mu(D),\ell_\mu(D))$, with
deterministic encoders and codes at every large blocklength, and it is
the nonnegative quadrant for $D\ge D_{\max}^\mu$. In particular the
rate--distortion function is $r_\mu(D)$.
\end{theorem}

The quantity $\theta$ is unique because
$\theta\mapsto Q\int\min\{\lambda,\theta\}\,d\mu$ is continuous and
strictly increasing until it reaches $D_{\max}^\mu$. The theorem applies
to stationary scalar sources, where $\mu$ is the distribution of the
values of the spectral density by Szeg\H{o}'s theorem
\cite{grenander1958}, and to the scalar locally stationary and
cyclostationary sources of \cite{part3}. Its interest is for sources
whose statistics vary too fast for a local spectrum and never repeat.

\begin{proof}
If $\Gamma_N$ is singular, $T^N$ lies in its range, the part of the
holder's observation orthogonal to the range is independent noise, and
the problem is that of the vector of coordinates in an orthonormal
basis of the range; the formulas below hold with the convention that a
zero eigenvalue costs nothing. \emph{Converse.} By
\cite[Lemma~24]{part3}, for every $\alpha$,
$\tfrac1N\VG(\alpha,ND;\Gamma_N,JI_N)=\alpha r_N+(1-\alpha)\ell_N$ with
$r_N=\int\tfrac12\log^+(\lambda/\theta_N)\,d\mu_N$, $\ell_N$ the
analogous integral, and $\theta_N$ the solution of
$Q\int\min\{\lambda,\theta_N\}\,d\mu_N=D$. The functions
$\theta\mapsto Q\int\min\{\lambda,\theta\}\,d\mu_N$ converge to the one
for $\mu$ uniformly on compact sets, because $\min\{\lambda,\theta\}$ is
bounded and Lipschitz on $[0,b]$, uniformly in $\theta$; with strict
monotonicity at $\theta$, $\theta_N\to\theta$. The integrands of $r$
and $\ell$ are bounded and jointly continuous in $(\lambda,\theta)$ on
$[0,b]\times[\theta/2,2\theta]$, so $r_N\to r_\mu(D)$ and
$\ell_N\to\ell_\mu(D)$. Every code of blocklength $N$ has
$\alpha R_N+(1-\alpha)L_N\ge\tfrac1N\VG(\alpha,ND_N;\Gamma_N,JI_N)$, and the
continuity of $r_\mu$ and $\ell_\mu$ in $D$ gives the converse as in the
proof of Theorem~\ref{thm:vector}.

\emph{Achievability.} The proof of \cite[Theorem~26]{part3} for the
bulk applies with two changes. The window is rotated into the
eigenbasis of $\Gamma_N$, which keeps the holder's noise white because
its gain is scalar, and the components are independent with variances
the eigenvalues. The bins are $[v_{i-1},v_i)$ with
$v_0=\epsilon_0\in(0,\theta)$, $v_i\le(1+\kappa)v_{i-1}$, $v_B>b$, and
$\mu(\{v_i\})=0$ for $0\le i<B$, and the components with eigenvalues
below $\epsilon_0$ are reproduced as zero. In the minimizing allocation
those components have $\delta=\lambda$, so they cost nothing there
either, and their distortion is the same in both. Each other bin is
coded with codes for white sources of variance $v_{i-1}$ and
\cite[Lemma~12]{part3}. Since $\mu_N\to\mu$ and the boundaries carry no
mass, the share of each bin tends to its $\mu$-mass, and the rest of
that proof gives every point of the quadrant. By
\cite[Lemma~24]{part3} the same allocation
$\delta=\min\{\lambda,\theta\}$ minimizes every weighted cost, so the
region is a quadrant. The construction works at every large $N$.
\end{proof}

\subsection{Sources driven by a uniquely ergodic map}

\begin{proposition}\label{prop:ergodic}
Let $\Theta$ be a compact metric space and $\tau:\Theta\to\Theta$ a
homeomorphism with exactly one invariant probability measure $\nu$.
Let $c:\Theta\times\{0,1,\dots,m\}\to\R$ be continuous in its first
argument, set $c(\vartheta,\Delta)=0$ for $\Delta>m$, and
for $\vartheta_0\in\Theta$ let $\Gamma_N(\vartheta_0)$ be the matrix with
entries $c(\tau^s\vartheta_0,t-s)$, $1\le s\le t\le N$, extended
symmetrically. If every $\Gamma_N(\vartheta_0)$ is positive semidefinite,
then $\|\Gamma_N(\vartheta_0)\|\le(2m+1)\max|c|$, and the eigenvalue
distributions $\mu_N(\vartheta_0)$ converge weakly to one probability
measure $\mu$, the same for every $\vartheta_0$, whose moments are
\[
\int\lambda^k\,d\mu=\int_\Theta g_k\,d\nu,\qquad
g_k(\vartheta)=\bigl(\Gamma(\vartheta)^k\bigr)_{00},
\]
with $\Gamma(\vartheta)$ the symmetric bi-infinite banded matrix with
entries $c(\tau^s\vartheta,t-s)$, $s\le t$, indexed by $\Z$. The moments of
$\mu_N(\vartheta_0)$ converge uniformly in $\vartheta_0$.
\end{proposition}

\begin{proof}
The bound on the norm is the Schur test. For $k\ge1$,
$\int\lambda^kd\mu_N=\tfrac1N\tr\Gamma_N^k=\tfrac1N\sum_s(\Gamma_N^k)_{ss}$,
and for $s$ at distance more than $km$ from both ends the diagonal
entry is a sum over paths of $k$ steps of length at most $m$, so it
equals $g_k(\tau^s\vartheta_0)$. The function $g_k$ is a finite sum of
products of continuous functions of $\vartheta$, hence continuous, and
each of the $2km$ other diagonal entries is bounded by
$((2m+1)\max|c|)^k$. By the uniform ergodic theorem for uniquely ergodic
maps \cite{oxtoby1952},
$\tfrac1N\sum_{s\le N}g_k(\tau^s\vartheta_0)\to\int g_kd\nu$ uniformly in
$\vartheta_0$. The measures $\mu_N$ live on a fixed compact interval,
where moments determine a measure and convergence of all moments
implies weak convergence.
\end{proof}

\begin{corollary}\label{cor:sampled}
Let $X_c(t)$, $t\in\R$, be a centered Gaussian process whose correlation
$c_c(t,\lambda)=\E[X_c(t)X_c(t+\lambda)]$ is continuous, periodic in $t$
with period $T_c$, and zero for $|\lambda|>\lambda_c$. Let $T_s>0$ with
$T_s/T_c$ irrational, let $\phi\in[0,T_c)$, and let
$T_i=X_c(iT_s+\phi)$, $i\ge1$. Then the eigenvalue distributions of the
window covariances converge to a measure $\mu$ that does not depend on
$\phi$, and with a scalar holder in white noise the region of
Theorem~\ref{thm:law} is the same for every sampling phase. In
particular the rate--distortion function of the sampled process exists
as a limit at every $D>0$, equals
$\int\tfrac12\log^+(\lambda/\theta)\,d\mu$, and does not depend on $\phi$.
\end{corollary}

\begin{proof}
Take $\Theta=\R/T_c\Z$, $\tau(\vartheta)=\vartheta+T_s$, which is uniquely
ergodic with Lebesgue measure because $T_s/T_c$ is irrational, and
$c(\vartheta,\Delta)=c_c(\vartheta,\Delta T_s)$, which vanishes for
$|\Delta|>\lambda_c/T_s$. Covariance matrices are positive semidefinite,
so Proposition~\ref{prop:ergodic} and Theorem~\ref{thm:law} apply.
\end{proof}

Tan, Dabora, and Poor \cite[Theorem~1]{tan2024} express the
rate--distortion function of this source, with a continuous correlation
and a finite correlation length, as an average over $\phi$ of limits
superior of the window programs, which by \cite[Lemma~24]{part3} are
$r_N$ in the notation of the proof of Theorem~\ref{thm:law}, and they
conjecture that each limit superior does not depend on $\phi$
\cite[Remark~1]{tan2024}. Corollary~\ref{cor:sampled} proves this, shows
that the limits exist, and gives their common value in closed form,
together with the leakage. Their average over the phase and their
maximum over it \cite[Corollary~1]{tan2024} are both equal to
$r_\mu(D)$. When $T_s/T_c$ is rational the map $\tau$ has finite orbits,
the moments of $\mu$ are averages of $g_k$ over the orbit of $\phi$
rather than over the circle, and they can depend on the orbit. For
example, let $X_c(t)=\sigma(t)Z(t)$ with $Z$ stationary, so that
$\Gamma_N=D_\sigma\Gamma^Z_ND_\sigma$ with $D_\sigma$ the diagonal matrix of
the sampled $\sigma$, and let $\sigma(t)^2=5+4\sin(2\pi t/T_c)$ and
$T_s=T_c/2$. The phase $\phi=0$ gives samples of variance $5$, and
$\phi=T_c/4$ gives variances alternating between $9$ and $1$, so
$\tfrac1N\log\det\Gamma_N$ differs between the two by $\log(5/3)$. At a
distortion small enough that every eigenvalue exceeds the water level,
the rate is $\tfrac1{2N}\log\det\Gamma_N-\tfrac12\log(D/Q)$, so the two
rate--distortion functions differ by $\tfrac12\log(5/3)$ bits there.

\subsection{Vector sources}

Theorem~\ref{thm:law} uses that a scalar source, a scalar weight, and a
scalar holder commute, so that the cost depends only on the
eigenvalues. Without that, the cost of a window is no longer a function
of its eigenvalues, but the region is still determined by averages of
block problems over the phase.

Let $\Theta$ be a compact metric space and $\tau$ a homeomorphism of
$\Theta$ such that every power $\tau^n$, $n\ge1$, has exactly one
invariant probability measure, the same $\nu$ for all $n$; an irrational
rotation of the circle is such a map. Let
$c:\Theta\times\{0,\dots,m\}\to\R^{p\times p}$ be continuous, with
$c(\vartheta,0)$ symmetric, and for $\vartheta_0\in\Theta$ let
$\Gamma_N(\vartheta_0)$ be the block matrix with the blocks
$c(\tau^s\vartheta_0,t-s)$ for $s\le t\le s+m$, their transposes below the
diagonal, and zero blocks otherwise. Assume a floor: there is $a>0$ with
$\Gamma_N(\vartheta_0)\succeq aI$ for every $N$ and $\vartheta_0$. Each
$\Gamma_N(\vartheta_0)$ is the leading principal submatrix of the next, so
they are the covariances of the windows of one centered Gaussian
process. The holder is memoryless, $S_t=HT_t+U_t$ with $U_t$ independent
$\mathcal N(0,\Sigma_U)$, and $J=H^{\mathsf T}\Sigma_U^{-1}H$. Write
$F_n(\vartheta)=\Gamma_n(\vartheta)$, let $b$ bound the norms of all
$\Gamma_N(\vartheta_0)$, which is finite by the Schur test, and let
$D_{\max}=\int\tr(Q\,c(\vartheta,0))\,d\nu$. For a probability measure
$\lambda$ on $\Theta$ let
\[
V_n(\lambda;\alpha,D)=\inf\Bigl\{\int\tfrac1n\VG\bigl(\alpha,nD(\vartheta);F_n(\vartheta),I_n\otimes J\bigr)\,d\lambda(\vartheta):\ D(\cdot)>0,\ \int D\,d\lambda\le D\Bigr\},
\]
the cost of describing independent blocks of $n$ symbols whose phases
have the law $\lambda$, under one budget, and let
$D_{\max}(\lambda)=\int\tfrac1n\tr((I_n\otimes Q)F_n(\vartheta))\,d\lambda$, the
distortion at which it vanishes. For $\lambda=\nu$, $D_{\max}(\nu)=D_{\max}$.

\begin{lemma}\label{lem:frozen}
For every $n$, $\alpha$, and $D>0$: (a) $V_n(\lambda;\alpha,\cdot)$ is
convex, nonincreasing, and continuous on $(0,\infty)$, and zero from
$D_{\max}(\lambda)$ on; (b) if $\lambda_k\to\lambda$ weakly then
$V_n(\lambda_k;\alpha,D)\to V_n(\lambda;\alpha,D)$; (c) $\lambda\mapsto V_n(\lambda;\alpha,D)$
is convex; (d) $|V_n(\lambda;\alpha,D)-V_n(\lambda;\alpha',D)|\le|\alpha-\alpha'|\,\tfrac p2\log(1+b\|J\|)$.
\end{lemma}

\begin{proof}
The floor makes $F_n$ invertible, so with $X=F_n^{-1/2}\Delta F_n^{-1/2}$
the frontier cost of each block is the program of \cite[Lemma~2]{part2}
at one point of the compact space $\Theta$, with the continuous data
$F_n^{1/2}(I_n\otimes Q)F_n^{1/2}$, $(I_n\otimes H)F_n^{1/2}$, and
$I_n\otimes\Sigma_U$. A field $X(\cdot)$ gives the allocation
$D(\vartheta)=\tfrac1n\tr(F_n^{1/2}(I_n\otimes Q)F_n^{1/2}X(\vartheta))$ with no
larger cost, and the continuous field of that lemma with one level
$\varpi_\lambda$ gives an allocation attaining its value, so $V_n(\lambda;\cdot)$
is the program of that lemma, up to the factor $1/n$. For
$D<D_{\max}(\lambda)$ that lemma gives (b) and the field with one level;
(a) holds because each frontier cost is convex in its budget and the
budget constraint is linear. For $D\ge D_{\max}(\lambda)$ the cost is zero,
and for $D=D_{\max}(\lambda)$, (b) follows from the continuity in (a) and
monotonicity, since $D_{\max}(\lambda_k)\to D_{\max}(\lambda)$. For (c), let
$g_\varpi(\vartheta)=\inf_{d>0}[\tfrac1n\VG(\alpha,nd;F_n(\vartheta),I_n\otimes J)+\varpi d]$.
For every feasible allocation and every $\varpi\ge0$,
$\int\tfrac1n\VG\,d\lambda\ge\int g_\varpi\,d\lambda-\varpi D$, and the field with
the level $\varpi_\lambda$ attains equality, with $\varpi_\lambda=0$ when
$D\ge D_{\max}(\lambda)$. So
$V_n(\lambda;\alpha,D)=\max_{\varpi\ge0}[\int g_\varpi\,d\lambda-\varpi D]$, a maximum
of affine functions of $\lambda$. For (d), each frontier cost is a minimum
over $\Delta$ of $\alpha r+(1-\alpha)\ell$, which is affine in $\alpha$ with
slope
$r-\ell=\tfrac12\log\det(I+F_n(I_n\otimes J))-\tfrac12\log\det(I+\Delta(I_n\otimes J))$,
which lies in $[0,\tfrac{pn}2\log(1+b\|J\|)]$ because
$0\prec\Delta\preceq F_n\preceq bI$; divided by $n$ this gives (d).
\end{proof}

\begin{theorem}\label{thm:ergodic}
For $0<D<D_{\max}$, $V_n(\nu;\alpha,D)$ converges as $n\to\infty$, uniformly
in $\alpha$, to a limit $V_\nu(\alpha,D)$, with
$|V_n(\nu;\alpha,D)-V_\nu(\alpha,D)|=O(1/n)$. For every $\vartheta_0$, $\mathcal A(D)$
is the set of pairs with $\alpha R+(1-\alpha)L\ge V_\nu(\alpha,D)$ for every
$\alpha$, the same for every starting phase, with deterministic
encoders, and the window cost
$\tfrac1N\VG(\alpha,ND;\Gamma_N(\vartheta_0),I_N\otimes J)$ tends to
$V_\nu(\alpha,D)$, uniformly in $\vartheta_0$.
\end{theorem}

The codes of the proof are built knowing $\vartheta_0$, as the codes of
\cite{tan2024} know the sampling phase; the region they achieve does not
depend on it.

\begin{proof}
Fix $\alpha$, $n$, and $D'$ with $D<D'<D_{\max}$.

\emph{Lower bound.} For an offset $s\in\{0,\dots,n-1\}$ split the window
into the first $s$ symbols, $k_s$ blocks of $n$, and a remainder of
fewer than $n$. Consecutive blocks are correlated only through the $m$
symbols on each side of the cut, so the quantity $\chi_b$ of
\cite[Lemma~10]{part3} is at most
$\chi=m^2\max_{\vartheta,\Delta}\|c(\vartheta,\Delta)\|_{\mathrm F}^2$. That lemma
applies with $a_0=a$ and $b_0=b$, to codes and also to the window
program, because its proof needs only the two informations, the Markov
chains of the description through the source, and finite conditional
entropies, which the optimal Gaussian test channel has
\cite[Sec.~II]{part3}. Dropping the first and the last pieces, whose
costs are nonnegative, it gives, for every code with distortion at most
$D'$,
\[
N[\alpha R+(1-\alpha)L]\ge\sum_{b}n\,\mathsf V_b(D_b)-c_1(k_s+2)\chi,
\]
where $b$ runs over the $k_s$ blocks, $\mathsf V_b(D_b)$ is the frontier
cost per symbol of block $b$ at its distortion $D_b$, and
$\sum_bnD_b\le ND'$. With the allocation that gives each block's phase
its $D_b$, and the average of the $D_b$ where two blocks share a phase,
which by convexity of the frontier cost in the budget does not raise the
sum, the right side is at least
$k_sn\,V_n(\mu_s;\alpha,ND'/(k_sn))-c_1(k_s+2)\chi$, with $\mu_s$ the uniform
measure on the phases of the blocks. Let $\underline k=\min_sk_s$,
$\omega_N=\underline kn/N$, and $D''_N=ND'/(\underline kn)$; both $\omega_N$
and $D''_N/D'$ tend to $1$. Since $V_n$ is nonincreasing in $D$ and
nonnegative, every offset gives
$\alpha R+(1-\alpha)L\ge\omega_NV_n(\mu_s;\alpha,D''_N)-c_1(N/n+2)\chi/N$.
Averaging over $s$ and using Lemma~\ref{lem:frozen}(c),
\[
\alpha R+(1-\alpha)L\ge\omega_NV_n(\bar\mu_N;\alpha,D''_N)-c_1(N/n+2)\chi/N,\qquad
\bar\mu_N=\tfrac1n\sum_s\mu_s .
\]
The measure $\bar\mu_N=\bar\mu_N(\vartheta_0)$ differs by $O(n/N)$ in total
variation from the empirical measure of $\tau^t\vartheta_0$, $t<N$, and by
the uniform ergodic theorem \cite{oxtoby1952} the integral of every
continuous function against it converges to its integral against $\nu$,
uniformly in $\vartheta_0$. We claim that
$V_n(\bar\mu_N(\vartheta_0);\alpha,D''_N)\to V_n(\nu;\alpha,D')$ uniformly in
$\vartheta_0$. If not, there are $N_k\to\infty$, $\vartheta_0^k$, and
$\epsilon>0$ with the two at least $\epsilon$ apart; but
$\bar\mu_{N_k}(\vartheta_0^k)\to\nu$ weakly, and for every $\eta>0$, eventually
$D'-\eta\le D''_{N_k}\le D'+\eta$, so by monotonicity and
Lemma~\ref{lem:frozen}(a),(b) the limit points lie between
$V_n(\nu;\alpha,D'+\eta)$ and $V_n(\nu;\alpha,D'-\eta)$, which tend to
$V_n(\nu;\alpha,D')$ as $\eta\to0$, a contradiction. Hence every pair
achievable at $D$ satisfies
$\alpha R+(1-\alpha)L\ge V_n(\nu;\alpha,D')-c_1\chi/n$, and so, letting
$D'\downarrow D$, $\alpha R+(1-\alpha)L\ge V_n(\nu;\alpha,D)-c_1\chi/n$. The
window cost satisfies the same bound, uniformly in $\vartheta_0$.

\emph{Upper bound.} Fix $\kappa\in(0,1)$. By uniform continuity and the
floor there is a finite partition of $\Theta$ into cells $E_i$ with
$\nu(\partial E_i)=0$ and points $\vartheta_i\in E_i$ such that
$\|F_n(\vartheta)-F_n(\vartheta_i)\|\le\kappa a$, hence
$(1-\kappa)F_n(\vartheta_i)\preceq F_n(\vartheta)$, for $\vartheta\in E_i$. Arrange
the window as blocks of $n$ symbols separated by guards of $m$ symbols.
Correlations have range $m$, so the blocks are independent, and block
$j$ has covariance $F_n$ at its phase. Because $\tau^{n+m}$ is uniquely
ergodic, the empirical measure of the phases of the blocks converges
weakly to $\nu$, uniformly in $\vartheta_0$, and since $\nu(\partial E_i)=0$,
the indicator of $E_i$ can be squeezed between continuous functions
whose integrals against $\nu$ differ by little, so the fraction of
blocks in $E_i$ tends to $\nu(E_i)$, uniformly in $\vartheta_0$. For each
cell with $\nu(E_i)>0$ take a code with a deterministic encoder for
independent copies of the Gaussian vector of $n$ symbols with
covariance $\mathsf S'_i=(1-\kappa)F_n(\vartheta_i)$, near a point of its
frontier at a distortion $d_i$ \cite[Theorem~3]{part1}, and apply it to
the blocks of that cell by \cite[Lemma~12]{part3}, whose proof does not
use that the groups are consecutive in time. Since
$F_n(\vartheta)-\mathsf S'_i\preceq\kappa(a+b)I$, the added distortion and
leakage per symbol are at most $\|Q\|p\kappa(a+b)$ and
$\|J\|p\kappa(a+b)/(2\ln2)$. The guards are reproduced as zero, at a
distortion of at most $c_gm/(n+m)$ per symbol with
$c_g=\max_\vartheta\tr(Qc(\vartheta,0))$, and so are the final partial
period, the blocks of cells of measure zero, and the fewer than one
group of blocks left over in each cell, whose share of the window tends
to zero. By \cite[Lemma~11(b)]{part2}, used with $\mathsf S'_i\preceq F_n(\vartheta_i)\preceq\mathsf S'_i/(1-\kappa)$,
the frontier cost of $\mathsf S'_i$ at the distortion $(1-\kappa)d$ is at
most that of $F_n(\vartheta_i)$ at $d$ plus $\tfrac p2\log\tfrac1{1-\kappa}$.
Let $\mathcal B$ be the set of pairs these codes achieve, over all
allocations $(d_i)$ with block budget $D-c_gm/(n+m)$. Each Gaussian-vector
region is jointly convex in the rate, the leakage, and the distortion
\cite[Theorem~5]{part1}, so, as in \cite[proof of Theorem~13]{part3},
$\mathcal B$ is closed, convex, and upward closed, and its support value
at $(\alpha,1-\alpha)$ is the program over the measure with masses
$\nu(E_i)$ at the points $\vartheta_i$ and the covariances $\mathsf S'_i$.
By the bounds above and Lemma~\ref{lem:frozen}(b),(d), as the cells
shrink and $\kappa\to0$ that value tends to
$V_n(\nu;\alpha,D-c_gm/(n+m))$ uniformly in $\alpha$. So every pair with
$\alpha R+(1-\alpha)L>V_n(\nu;\alpha,D-c_gm/(n+m))+\gamma$ for every $\alpha$, for
some $\gamma>0$, is achievable at $D$, with codes at every large $N$,
uniformly in $\vartheta_0$, and the window cost is at most
$V_n(\nu;\alpha,D-c_gm/(n+m))+o(1)$.

\emph{The limit and the region.} The functions $V_n(\nu;\alpha,\cdot)$ are
convex and nonincreasing, and at $D/4$ they are at most
$V_n(\nu;1,D/4)$, because $\ell\le r$, which is at most the rate of a
Gaussian vector with covariance at most $bI$ at that distortion,
bounded in $n$. So on $[D/2,D]$ their slopes are bounded by a constant
$L_D$ that depends on neither $n$ nor $\alpha$. With $\epsilon_n=c_gm/(n+m)$,
the two bounds give, for every achievable pair and for every
$\alpha$,
\[
V_n(\nu;\alpha,D)-\frac{c_1\chi}n\le\alpha R+(1-\alpha)L,
\]
while every pair above $V_n(\nu;\alpha,D)+L_D\epsilon_n$ for every $\alpha$ is
in $\mathcal A(D)$. The set of pairs above
$V_{n'}(\nu;\alpha,D)+L_D\epsilon_{n'}$ for every $\alpha$ has, at
$(\alpha,1-\alpha)$, the support value $V_{n'}(\nu;\alpha,D)+L_D\epsilon_{n'}$,
attained by the pair of rates and leakages of an optimal allocation of
the program at $\alpha$, which lies in the set because that program is a
minimum at every other $\alpha$. Since that set lies in $\mathcal A(D)$,
the first bound gives
$V_n(\nu;\alpha,D)-c_1\chi/n\le V_{n'}(\nu;\alpha,D)+L_D\epsilon_{n'}$ for all $n$
and $n'$, so the sequence $V_n(\nu;\alpha,D)$ is Cauchy uniformly in
$\alpha$, with limit $V_\nu(\alpha,D)$ and error at most
$\max\{c_1\chi/n,L_D\epsilon_n\}$. The first bound shows
that every achievable pair satisfies $\alpha R+(1-\alpha)L\ge V_\nu(\alpha,D)$.
Conversely, if $(R,L)$ satisfies this for every $\alpha$, then for every
$\gamma>0$ the pair $(R+\gamma,L+\gamma)$ lies above $V_n(\nu;\alpha,D)+L_D\epsilon_n$
for every $\alpha$ once $n$ is large, so it is achievable, and $(R,L)$ is
in the closure $\mathcal A(D)$. The window cost lies between the two
bounds, uniformly in $\vartheta_0$. None of the bounds involves
$\vartheta_0$.
\end{proof}

For a scalar source and a scalar holder, Theorems~\ref{thm:law}
and~\ref{thm:ergodic} describe the same region where both apply, so
$V_\nu$ is the closed form of Theorem~\ref{thm:law}, which needs neither
the floor nor total unique ergodicity. For a vector source we do not
have a closed form; $V_n(\nu;\cdot)$ is a program over fields on $\Theta$
with one level, and the proof bounds its error by $O(1/n)$, with
constants that depend on the floor $a$ and on $D$. No rate is claimed
for the window cost, whose ergodic averages can converge slowly.

\section{Open Problems}
\label{sec:open}

\textbf{Holders with infinite memory.} Theorem~\ref{thm:vector} is for a
holder whose filter has finitely many taps. Truncating a filter with
infinitely many, as in \cite[Theorem~22]{part2}, does not work directly
for an unstable source, because the tail of the filter acts on a signal
that grows. For a stable rational filter a joint state-space
realization of the source and the filter is the natural route, with the
hidden poles those cancelled by the filter's zeros; its information
matrix on a window is no longer a finite-rank perturbation of a block
Toeplitz one.

\textbf{Holders with memory, and closed forms, under fast variation.}
Theorem~\ref{thm:ergodic} is for a memoryless holder and needs a floor
on the covariances, which Theorem~\ref{thm:law} does not. For a vector
source its limit is a limit of block programs, not a closed form. For a
rotation it should equal the limit of the regions of the
cyclostationary sources with rational rotation numbers, which
\cite[Theorem~18]{part3} computes.

\textbf{Poles on the unit circle.} For a random walk the bridge of
Lemma~\ref{lem:bridge} is a discrete Brownian bridge, whose variance
grows with $N$, so the growth is no longer carried by the final state
alone. A pole of modulus one contributes $\log1=0$ to $\rho_\star$, and we
expect the region to be $\mathcal R_f(D)$ without a shift, with $f$ now
unbounded and $D_{\max}^f=\infty$; the rate part is classical
\cite{gray1970,grayhashimoto2008}.

\textbf{Processes that are not Gaussian.} Linear processes with
infinite dependence and innovations that are not Gaussian remain open,
as in \cite{part3}.

\appendices

\section{Numerical Checks}
\label{app:numver}
\sloppy

The statements were checked numerically, independently of the proofs.
The scripts and their recorded outputs are in the public repository
\texttt{https://github.com/ahb-sjsu/geometric-observation}, directory
\texttt{paper/tit-unstable-ergodic}; a tag will fix the version that goes
with this draft. They use the solver \texttt{spectral\_core.py} of
\cite{part2}. No proof depends on them. The pass criteria of the checks
of Section~\ref{sec:law} were written into the script before its first
run. Those of Section~\ref{sec:vector} were written after exploratory
runs of the same computations on the same sources, and are therefore a
record of agreement, not a test fixed in advance.

\texttt{verify\_unstable\_vector.py} uses twelve sources: first-order
sources in $\R^2$ with a diagonal matrix, a rotation (a complex pair of
modulus $1.2$), and a Jordan block (eigenvalue $1.3$), each with a
memoryless holder that sees the unstable mode and one that does not,
and with a triangular matrix with unstable eigenvalues $1.4$ and $1.2$
and a holder from which only the first is hidden; a
second-order source in $\R^2$ with a pole at $2$ hidden from the holder
next to a visible pole at $-1.19$, built with dyadic coefficients so that
the pole is hidden exactly, and the same source with a holder that sees
both; the scalar source with pole $1.5$ and the holder filters
$1-1.5z$, which hides the pole, and $1-0.5z$, which does not; and one
source described below. The covariances of the final state are
computed with $320$-digit arithmetic, because at $N=400$ they span more
than $160$ decimal orders. The difference in Lemma~\ref{lem:bridge}
vanishes outside the last $d$ blocks exactly and has rank at most $pd$,
in every case at $N=12$. The precision of a stationary window used in
Lemma~\ref{lem:stationary} differs from $T_N(f^{-1})$ by a positive
semidefinite matrix of rank at most $2pd$ on the first and last $d$
blocks, to $10^{-12}$, at $N=16$. In every source but the twelfth, the
identity of Lemma~\ref{lem:innovations} holds to $2\times10^{-14}$, and the
growth rates $\tfrac1{2N}\ln\det\Cov(Z)$ and
$\tfrac1{2N}\ln\det\Cov(Z\mid S^N)$ approach $\rho_\star\ln2$ and $\rhid\ln2$
with $N$ times the error below $2$ at $N=400$
(Lemma~\ref{lem:endpoint}(a),(b)), and the smallest eigenvalue of the
posterior covariance stays above the bound of
Lemma~\ref{lem:endpoint}(c), by a factor of at least $1.7$, at every $N$.
For the window cost of Theorem~\ref{thm:vector}, on the first-order
source with a visible and with a hidden pole and on the second-order
source with the hidden pole, at three weights and $D=0.3D_{\max}^f$, the
cost minus its predicted limit shrinks by a factor of $8$ from $N=40$ to
$N=320$, where it is at most $2\times10^{-3}$ bits. The same holds for the
scalar source with the holder that sees the innovations, at most
$9\times10^{-4}$ bits at $N=320$; that run was made before
Theorem~\ref{thm:vector} covered holders with memory and has no pass
criterion. The window cost is computed by capping the exploding part of
$\Gamma_N$ at a level $M$, which gives an upper bound that decreases to
the cost as $M$ grows. The criterion that the cost change by at most
$10^{-6}$ bits between $M=10^3$ and $M=10^5$ failed for the second-order
source at $N=40$ and $80$, by up to $1.9\times10^{-6}$ bits. A follow-up
with $M=10^7$ changed those costs by at most $2\times10^{-8}$ bits more,
so the values are settled to that accuracy, three orders of magnitude
below the gaps they measure. The first run is recorded with its
failure.

The twelfth source shows the discontinuity of $\rhid$. Built in double
precision with a pole at $1.6$ meant to be hidden, its mode reaches the
holder with gain about $6\times10^{-17}$. The posterior growth rate
$\tfrac1{2N}\ln\det\Cov(Z\mid S^N)$ is $0.40$ and $0.44$ at $N=25$ and
$50$, as for a hidden pole, and then falls to $0.40$, $0.20$, and $0.10$
at $N=100$, $200$, and $400$, as the holder begins to see the mode; the
innovation identity computed with the pole counted as hidden fails by
$0.94=2\ln1.6$, the amount of one visible pole.

\texttt{verify\_spectral\_law.py} checks Corollary~\ref{cor:sampled} on
the process $\sigma(t)Z(t)$ with $\sigma(t)^2=5+\tfrac13\sin2\pi t$, the
variance profile of the example of \cite{tan2024}, and $Z$ stationary
with a triangular correlation of length $2.5$ samples, sampled at
$T_s=1/(3+\varepsilon)$ with $\varepsilon=(\sqrt5-1)/2$. Over four sampling
phases, the largest difference in the rate or the leakage of
Theorem~\ref{thm:law}, at two distortions, falls from
$1.5\times10^{-4}$ bits at $N=500$ to $7.3\times10^{-6}$ bits at $N=8000$.
The first six moments of the eigenvalue distribution at $N=8000$ agree
with the formula of Proposition~\ref{prop:ergodic} to
$2.4\times10^{-4}$ relative. Two rational sampling rates, with periods of
$7$ and $3$ samples, also show no measurable phase dependence for this
source, whose modulation is weak; the second was added after the first
run. The exact example after Corollary~\ref{cor:sampled} shows phase
dependence at a rational rate.

\texttt{verify\_vector\_ergodic.py} checks Theorem~\ref{thm:ergodic} on a
source in $\R^2$ driven by the same rotation: a stationary vector
moving average of order one, multiplied by a continuous matrix of the
phase, plus a white floor of variance $0.2$, with the holder
$H=(1,0.5)$, at three weights and $D=0.3D_{\max}$. Its first run, with
criteria fixed before it, failed all six checks: the spread of the
window cost over four phases at $N=200$ was $3.9\times10^{-3}$ bits
against a limit of $2\times10^{-3}$, and the frozen-block cost at $n=16$
was $1.1\times10^{-2}$ bits above the window cost at $N=200$ against a
limit of $5\times10^{-3}$. The relative parts of both criteria passed, and
both sequences were decreasing with differences that halved, as the
proof predicts; the absolute tolerances were too tight for those sizes.
The record of that run is kept. A follow-up to $n=64$ and $N=400$, with
criteria written before it ran, passed all nine checks: both sequences
decrease, their extrapolations $2V_{64}-V_{32}$ and $2W_{400}-W_{200}$ agree
to $1.6\times10^{-5}$ bits or better (at $\alpha=0$, $0.84456$ and $0.84457$),
and the phase spread at $N=400$ is at most a quarter of that at $N=100$.

Not checked numerically: the coding theorems, Lemmas~\ref{lem:perturb},
\ref{lem:circulant}, \ref{lem:groups}, \ref{lem:quantizer},
and~\ref{lem:frozen}, which
the window-cost check exercises only through its result, and
Proposition~\ref{prop:ergodic} for maps other than a rotation.
Table~\ref{tab:checks} lists what each check covers.

\begin{table}[!t]
\centering
\footnotesize
\caption{What the numerical checks cover. A check of a quantity is not a
check of the coding theorem that uses it.}
\label{tab:checks}
\begin{tabular}{lll}
\hline
Statement & Check & What is compared\\
\hline
Lemma~\ref{lem:bridge} & \texttt{bridge} & window precision against $T_N(f^{-1})$\\
Lemma~\ref{lem:stationary} & \texttt{stationary} & stationary window precision against $T_N(f^{-1})$\\
Lemma~\ref{lem:innovations} & \texttt{innovations} & Kalman innovations against the spectral integral\\
Lemma~\ref{lem:endpoint} & \texttt{prior\_growth}, \texttt{posterior\_growth}, \texttt{posterior\_floor} & covariances of the final state\\
Theorem~\ref{thm:vector} & \texttt{window\_cost} & window cost against $\alpha\rho_\star+(1-\alpha)\rhid+V_f$\\
Proposition~\ref{prop:ergodic} & \texttt{moments} & eigenvalue moments against the ergodic average\\
Corollary~\ref{cor:sampled} & \texttt{phase\_independence} & rate and leakage over sampling phases\\
Theorem~\ref{thm:ergodic} & \texttt{monotone}, \texttt{common\_limit}, \texttt{phase\_spread} & frozen blocks against window costs\\
\hline
\end{tabular}
\end{table}

\section*{Acknowledgment}

In accordance with the IEEE policy on AI-generated text, the author
discloses the use of Anthropic's Claude models (Claude Fable 5.2 and
Claude Opus 5.5) through the Claude Code environment in 2026, for
drafting prose and proofs, for the prior-art search, for the check
scripts, and for checking the proofs. The author is responsible for the
content.

\begin{thebibliography}{99}
\bibitem{part1} A. H. Bond, ``Rate, leakage, and distortion for Gaussian
sources: Water-filling on a tilted weight,'' submitted to \emph{IEEE
Trans. Inf. Theory}, 2026, archived at
\url{https://doi.org/10.5281/zenodo.23127802}.
\bibitem{part2} A. H. Bond, ``Rate, leakage, and distortion for
stationary Gaussian sources: Water-filling on a tilted spectrum,''
manuscript, 2026, version \texttt{tit-stationary-1.2} at
\url{https://github.com/ahb-sjsu/geometric-observation/tree/tit-stationary-1.2/paper/tit-stationary}.
\bibitem{part3} A. H. Bond, ``Rate, leakage, and distortion beyond
stationary Gaussian sources: Water-filling over time and frequency,''
manuscript, 2026, version \texttt{tit-beyond-stationary-1.3} at
\url{https://github.com/ahb-sjsu/geometric-observation/tree/tit-beyond-stationary-1.3/paper/tit-beyond-stationary}.
\bibitem{gray1970} R. M. Gray, ``Information rates of autoregressive
processes,'' \emph{IEEE Trans. Inf. Theory}, vol.~16, no.~4,
pp.~412--421, Jul. 1970.
% abstract and metadata read; full text not read.
\bibitem{hashimoto1980} T. Hashimoto and S. Arimoto, ``On the
rate-distortion function for the nonstationary Gaussian autoregressive
process,'' \emph{IEEE Trans. Inf. Theory}, vol.~26, no.~4,
pp.~478--480, Jul. 1980.
% metadata and its restatement in grayhashimoto2008 read; full text not read.
\bibitem{grayhashimoto2008} R. M. Gray and T. Hashimoto, ``A note on
rate-distortion functions for nonstationary Gaussian autoregressive
processes,'' \emph{IEEE Trans. Inf. Theory}, vol.~54, no.~3,
pp.~1319--1322, Mar. 2008.
% preprint (arXiv cs/0702150) read by the literature agent.
\bibitem{sahai2007} A. Sahai and S. Mitter, ``Source coding and channel
requirements for unstable processes,'' arXiv:cs/0610143, 2007.
% arXiv v2 read by the literature agent; journal publication not verified.
\bibitem{derpich2021} M. S. Derpich, M. M\"uller, and J. \O stergaard,
``The entropy gain of linear systems and some of its implications,''
\emph{Entropy}, vol.~23, no.~8, Art. no.~947, 2021.
% read through a summary by the literature agent. verify on final.
\bibitem{gutierrez2018} J. Guti\'errez-Guti\'errez,
M. Z\'arraga-Rodr\'iguez, P. M. Crespo, and X. Insausti, ``Rate
distortion function of Gaussian asymptotically WSS vector processes,''
\emph{Entropy}, vol.~20, no.~9, Art. no.~719, 2018.
\bibitem{gutierrez2024} J. Guti\'errez-Guti\'errez,
J. Barasoain-Echepare, M. Z\'arraga-Rodr\'iguez, and X. Insausti,
``Computation of the fundamental limits of data compression for certain
nonstationary ARMA vector sources,'' \emph{Probl. Inf. Transm.},
vol.~60, no.~1, pp.~1--11, 2024.
% not read here. Its Theorem 4 assumes companion spectral radius below one: reported on 2026-11-16 by a reader who accessed the article. verify on final.
\bibitem{toms1971} W. Toms and T. Berger, ``Information rates of
stochastically driven dynamic systems,'' \emph{IEEE Trans. Inf.
Theory}, vol.~17, no.~1, pp.~113--114, Jan. 1971.
% correspondence item, DOI 10.1109/TIT.1971.1054569. Abstract and index records only; full text not read. Author list differs between records; verify on IEEE Xplore. Whether unstable systems are allowed is unknown.
\bibitem{braslavsky1999} J. H. Braslavsky, M. M. Seron, D. Q. Mayne,
and P. V. Kokotovi\'c, ``Limiting performance of optimal linear
filters,'' \emph{Automatica}, vol.~35, no.~2, pp.~189--199, 1999.
% technical report version read by the literature agent.
\bibitem{abakasanga2020} E. Abakasanga, N. Shlezinger, and R. Dabora,
``On the rate-distortion function of sampled cyclostationary Gaussian
processes,'' \emph{Entropy}, vol.~22, no.~3, Art. no.~345, 2020.
% arXiv 2003.00262 read by the literature agent.
\bibitem{tan2024} Z. Tan, R. Dabora, and H. V. Poor, ``On the
rate-distortion function for sampled cyclostationary Gaussian processes
with memory,'' in \emph{Proc. IEEE Int. Symp. Inf. Theory (ISIT)}, 2024;
extended version arXiv:2405.11405.
% arXiv v2 pages 1-6 read (Theorem 1, Corollary 1, Remark 1). verify the ISIT pages on final.
\bibitem{andersonmoore1979} B. D. O. Anderson and J. B. Moore,
\emph{Optimal Filtering}. Englewood Cliffs, NJ: Prentice-Hall, 1979.
% verify the chapter cited for convergence of the Riccati recursion to the stabilizing solution under detectability and stabilizability, on final.
\bibitem{kailath2000} T. Kailath, A. H. Sayed, and B. Hassibi,
\emph{Linear Estimation}. Upper Saddle River, NJ: Prentice Hall, 2000.
% cited for the same convergence result; chapter not checked. verify on final.
\bibitem{rosenblum1985} M. Rosenblum and J. Rovnyak, \emph{Hardy
Classes and Operator Theory}. New York, NY: Oxford Univ. Press, 1985.
% cited for the matrix Fejer-Riesz theorem. verify on final.
\bibitem{oxtoby1952} J. C. Oxtoby, ``Ergodic sets,'' \emph{Bull. Amer.
Math. Soc.}, vol.~58, pp.~116--136, 1952.
% citation verified by the literature agent; text not read.
\bibitem{grenander1958} U. Grenander and G. Szeg\H{o}, \emph{Toeplitz
Forms and Their Applications}. Berkeley, CA: Univ. of California Press,
1958.
\end{thebibliography}

\end{document}
